\documentclass[UTF8,11pt,a4paper]{ctexart}
\usepackage{amsmath,amssymb,geometry,enumitem,hyperref,array,longtable,booktabs}
\geometry{margin=2.1cm}\setlist{itemsep=2pt,topsep=3pt}
\hypersetup{colorlinks=true,linkcolor=blue}
\setlength{\emergencystretch}{2em}
\ctexset{section={beforeskip=1.8ex,afterskip=0.8ex},subsection={beforeskip=1.2ex,afterskip=0.5ex}}
\setlength{\abovedisplayskip}{3pt}
\setlength{\belowdisplayskip}{3pt}
\setlength{\abovedisplayshortskip}{1pt}
\setlength{\belowdisplayshortskip}{1pt}
\newcommand{\APS}[2]{\par\noindent\textbf{中文。} #1\par\noindent\textbf{English.} #2\par\smallskip}
\newcommand{\EN}[1]{\par\noindent \textbf{English.} #1\par}
\newcommand{\FORMULA}[2]{\par\noindent\textbf{公式名称 / Formula name.} #1 / #2\par}
\newcommand{\READ}[1]{\par\noindent \textbf{公式英文读法 / How to read it.} #1\par}
\newcommand{\EX}[2]{\par\noindent \textbf{例子 / Example.} #1\par\EN{#2}}
\newcommand{\ASK}[4]{\par\noindent \textbf{可能问题 / Possible question.} #1\par\EN{#2}\noindent \textbf{回答要点 / Key answer.} #3\par\EN{#4}\medskip}
\newcommand{\dd}{\mathrm{d}}
\newcommand{\RR}{\mathbb{R}}
\newcommand{\br}{\mathbf{r}}
\newcommand{\bn}{\mathbf{n}}
\newcommand{\bT}{\mathbf{T}}
\newcommand{\bN}{\mathbf{N}}
\newcommand{\bB}{\mathbf{B}}
\newcommand{\be}{\mathbf{e}}
\newcommand{\bu}{\mathbf{u}}
\newcommand{\bv}{\mathbf{v}}
\newcommand{\bw}{\mathbf{w}}
\newcommand{\ba}{\mathbf{a}}
\newcommand{\bb}{\mathbf{b}}
\newcommand{\bx}{\mathbf{x}}
\newcommand{\bs}{\mathbf{s}}
\newcommand{\bc}{\mathbf{c}}
\newcommand{\bz}{\mathbf{z}}
\newcommand{\bgamma}{\boldsymbol{\gamma}}
\newcommand{\balpha}{\boldsymbol{\alpha}}
\newcommand{\bbeta}{\boldsymbol{\beta}}
\newcommand{\bomega}{\boldsymbol{\omega}}
\title{微分几何 APS 面试准备\\[4pt]\large 第 1--6 章：曲线论、曲面论与测地几何}
\author{Differential Geometry: Bilingual Review}
\date{}
\pagestyle{plain}
\begin{document}\maketitle

\section*{课程概述 | Course overview}
\APS{微分几何用微积分研究曲线与曲面的局部和整体性质。本讲义的主线是：先用标架和向量函数把几何对象写成可微映射；再在弧长参数下用曲率和挠率完全刻画空间曲线（Frenet 标架与曲线论基本定理）；然后把曲面放在三维欧氏空间中，用第一基本形式度量长度、角度和面积，用第二基本形式描述弯曲；接着由自然标架的运动公式得到 Gauss--Codazzi 方程，并证明 Gauss 绝妙定理：Gauss 曲率只由第一基本形式决定；最后研究曲面上的测地线、平行移动与 Gauss--Bonnet 公式，把局部曲率与整体拓扑联系起来。}{Differential geometry uses calculus to study the local and global properties of curves and surfaces. The main line of this review is: first write geometric objects as differentiable maps using frames and vector functions; then, in arc-length parametrization, characterize a space curve completely by its curvature and torsion (the Frenet frame and the fundamental theorem of curve theory); next place a surface in three-dimensional Euclidean space, measure length, angle and area with the first fundamental form, and describe bending with the second fundamental form; then obtain the Gauss--Codazzi equations from the motion formulas of the natural frame and prove Gauss's Theorema Egregium: the Gaussian curvature is determined by the first fundamental form alone; finally study geodesics, parallel transport and the Gauss--Bonnet formula, which links local curvature to global topology.}
\APS{这里有三个层次的对比。局部与整体：曲率和挠率是局部量，Gauss--Bonnet 把曲率积分与 Euler 示性数这样的整体不变量联系起来。外蕴与内蕴：第二基本形式依赖曲面在 $\RR^3$ 中的嵌入方式，第一基本形式只依赖曲面上的度量；绝妙定理说明 Gauss 曲率其实是内蕴量，这是整个课程最重要的观念转折。计算与结构：面试时既要能算 $\kappa,\tau,E,F,G,L,M,N,K,H$，也要能说清这些量为什么这样定义、满足什么方程。}{Three contrasts organize the subject. Local versus global: curvature and torsion are local quantities, while Gauss--Bonnet links an integral of curvature to a global invariant, the Euler characteristic. Extrinsic versus intrinsic: the second fundamental form depends on how the surface sits in $\RR^3$, whereas the first fundamental form depends only on the metric on the surface; the Theorema Egregium says that the Gaussian curvature is in fact intrinsic, and this is the key conceptual turn of the whole course. Computation versus structure: in an interview one must both compute $\kappa,\tau,E,F,G,L,M,N,K,H$ and explain why these quantities are defined as they are and which equations they satisfy.}
\APS{两个导入例子概括了全书。第一个是螺旋线：它的曲率和挠率都是常数，说明一条非平面曲线可以只靠两个数被完全确定；第二个是球面：第一、第二基本形式都只有对角项，算出的 Gauss 曲率为正常数 $1/R^2$，说明"球面不能摊平"这一直觉可以被一个数精确表达。}{Two opening examples summarize the book. The first is the helix: its curvature and torsion are both constant, showing that a non-planar curve can be determined completely by two numbers. The second is the sphere: both fundamental forms are diagonal, and the resulting Gaussian curvature is the positive constant $1/R^2$, showing that the intuition that a sphere cannot be flattened is captured precisely by a single number.}
\EX{螺旋线 $\bgamma(t)=(a\cos t,a\sin t,bt)$。直接计算得速度 $|\bgamma'|=\sqrt{a^2+b^2}$、曲率 $\kappa=a/(a^2+b^2)$、挠率 $\tau=b/(a^2+b^2)$。取 $a=b=1$ 得 $\kappa=\tau=1/2$；取 $b=0$ 退化为平面圆，$\kappa=1/a$、$\tau=0$。曲率度量曲线偏离直线的程度，挠率度量曲线离开密切平面的程度。}{Helix $\bgamma(t)=(a\cos t,a\sin t,bt)$. A direct computation gives speed $|\bgamma'|=\sqrt{a^2+b^2}$, curvature $\kappa=a/(a^2+b^2)$ and torsion $\tau=b/(a^2+b^2)$. For $a=b=1$ one gets $\kappa=\tau=1/2$; for $b=0$ the curve degenerates to a plane circle with $\kappa=1/a$ and $\tau=0$. Curvature measures how far the curve deviates from a straight line, and torsion measures how far it leaves its osculating plane.}
\EX{球面 $\br(u,v)=R(\sin u\cos v,\sin u\sin v,\cos u)$，其中 $0<u<\pi$、$0<v<2\pi$。计算得
\[
E=R^2,\quad F=0,\quad G=R^2\sin^2u;\qquad L=-R,\quad M=0,\quad N=-R\sin^2u,
\]
故 $K=(LN-M^2)/(EG-F^2)=1/R^2$、$H=(EN-2FM+GL)/(2(EG-F^2))=-1/R$（取外法向）。这两组系数就是本讲义第 3、4 章要系统研究的对象。}{Sphere $\br(u,v)=R(\sin u\cos v,\sin u\sin v,\cos u)$ with $0<u<\pi$ and $0<v<2\pi$. A direct computation gives $E=R^2$, $F=0$, $G=R^2\sin^2u$ and $L=-R$, $M=0$, $N=-R\sin^2u$, hence $K=(LN-M^2)/(EG-F^2)=1/R^2$ and $H=(EN-2FM+GL)/(2(EG-F^2))=-1/R$ for the outward normal. These two sets of coefficients are exactly what Chapters 3 and 4 study systematically.}
\subsection*{主线地图 | Road map}
\begin{center}\small
\begin{tabular}{@{}>{\raggedright\arraybackslash}p{1.4cm}>{\raggedright\arraybackslash}p{4.8cm}>{\raggedright\arraybackslash}p{9.2cm}@{}}
\toprule
章 Chapter & 核心对象 Object & 关键结论 Key conclusion \\
\midrule
1 & 标架、向量函数 frame, vector function & 把曲线曲面写成可微映射，用活动标架求导 \\
2 & 空间曲线 space curve & Frenet 公式；$\kappa,\tau$ 在刚体运动下完全决定曲线 \\
3 & 曲面、第一基本形式 first fundamental form & $E,F,G$ 给出长度、角度、面积；可展曲面 \\
4 & 第二基本形式 second fundamental form & $L,M,N$；主曲率、$K,H$；Dupin 标形 \\
5 & 基本方程 fundamental equations & Gauss--Codazzi 方程；绝妙定理 $K=K(E,F,G)$ \\
6 & 测地几何 geodesic geometry & 测地线、平行移动、Gauss--Bonnet 公式 \\
\bottomrule
\end{tabular}
\end{center}

\section[预备知识]{预备知识：标架与向量函数 | Preliminaries: frames and vector functions}
\subsection{三维欧氏空间中的标架 | Frames in Euclidean 3-space}
\APS{$\RR^3$ 中的一个标架由一个原点 $O$ 和三个右手正交单位向量 $\be_1,\be_2,\be_3$ 组成，即 $\be_i\cdot\be_j=\delta_{ij}$、$\be_1\times\be_2=\be_3$。任何向量 $\ba$ 在这组基下写成 $\ba=a_1\be_1+a_2\be_2+a_3\be_3$，坐标由 $a_i=\ba\cdot\be_i$ 给出。内积 $\ba\cdot\bb=\sum a_ib_i$、外积 $\ba\times\bb$ 的坐标也可用行列式写出。标架的作用是：把"绝对"的几何对象换成相对于一组方向的分量，从而把几何问题化为数组的运算。}{A frame in $\RR^3$ consists of an origin $O$ and three right-handed orthonormal vectors $\be_1,\be_2,\be_3$, that is $\be_i\cdot\be_j=\delta_{ij}$ and $\be_1\times\be_2=\be_3$. Every vector $\ba$ is written as $\ba=a_1\be_1+a_2\be_2+a_3\be_3$, with coordinates $a_i=\ba\cdot\be_i$. The inner product is $\ba\cdot\bb=\sum a_ib_i$, and the cross product $\ba\times\bb$ also has a determinant formula in coordinates. The role of a frame is to replace an absolute geometric object by components relative to a chosen set of directions, turning geometry into array computations.}
\APS{若标架随参数 $t$ 变化，就得到活动标架 $\{p(t);\be_1(t),\be_2(t),\be_3(t)\}$。每个 $\be_i'$ 都是标架向量的线性组合，写成 $\be_i'=\sum_j\omega_{ij}\be_j$。由正交性 $(\be_i\cdot\be_j)'=0$ 得 $\omega_{ij}+\omega_{ji}=0$，所以 $3\times3$ 系数矩阵 $\Omega=(\omega_{ij})$ 反对称，只有三个独立分量。在三维中反对称矩阵总可以写成一个向量 $\bomega$ 与向量的叉积，即 $\be_i'=\bomega\times\be_i$，$\bomega$ 称为标架的角速度向量。}{If the frame varies with a parameter $t$, one obtains a moving frame $\{p(t);\be_1(t),\be_2(t),\be_3(t)\}$. Each derivative $\be_i'$ is a linear combination of the frame vectors, written $\be_i'=\sum_j\omega_{ij}\be_j$. Orthonormality, $(\be_i\cdot\be_j)'=0$, gives $\omega_{ij}+\omega_{ji}=0$, so the $3\times3$ coefficient matrix $\Omega=(\omega_{ij})$ is skew-symmetric and has only three independent entries. In three dimensions every skew-symmetric matrix can be written as the cross product with a vector $\bomega$, that is $\be_i'=\bomega\times\be_i$, and $\bomega$ is called the angular velocity vector of the frame.}
\APS{这个"标架导数 = 反对称矩阵乘以标架"的格式是本课程反复出现的母题。Frenet 公式是它沿曲线的特例，曲面的自然标架运动公式是它沿曲面的二维版本；第 7 章用外微分把它写成结构方程 $\dd\be_i=\sum_j\omega_{ij}\be_j$。面试时可以用一句话概括活动标架法：不研究曲线或曲面本身，而研究附着在它上面的标架如何转动，曲率就是转动的速率。}{The pattern "derivative of a frame = a skew-symmetric matrix times the frame" is the recurring motif of the course. The Frenet formulas are its one-dimensional version along a curve, and the motion formulas of the natural frame are its two-dimensional version on a surface; Chapter 7 writes it as the structure equations $\dd\be_i=\sum_j\omega_{ij}\be_j$ in exterior calculus. In one sentence, the moving-frame method studies not the curve or surface itself but how a frame attached to it rotates, and curvature is the rate of that rotation.}

\subsection{向量函数 | Vector functions}
\APS{向量函数 $\br(t)=(x(t),y(t),z(t))$ 是三个数量函数组成的三维向量。极限、连续、导数与积分都逐分量定义：$\br'(t)=(x'(t),y'(t),z'(t))$，$\int_a^b\br(t)\dd t=(\int_a^bx,\int_a^by,\int_a^bz)$。因此一元微积分的定理逐条成立：可导必连续，$\br(t)-\br(t_0)=\int_{t_0}^t\br'(u)\dd u$，以及求导法则
\[
(\br\cdot\bs)'=\br'\cdot\bs+\br\cdot\bs',\qquad
(\br\times\bs)'=\br'\times\bs+\br\times\bs',\qquad
(\br\circ\varphi)'=\varphi'\,\br'(\varphi).
\]
若 $\br$ 有 $k$ 阶连续导数，记 $\br\in C^k$。}{A vector function $\br(t)=(x(t),y(t),z(t))$ is a three-dimensional vector made of three scalar functions. Limits, continuity, derivatives and integrals are defined componentwise: $\br'(t)=(x'(t),y'(t),z'(t))$ and $\int_a^b\br(t)\dd t=(\int_a^bx,\int_a^by,\int_a^bz)$. One-variable calculus therefore carries over theorem by theorem: differentiability implies continuity, $\br(t)-\br(t_0)=\int_{t_0}^t\br'(u)\dd u$, and the product rules displayed, where $\circ$ denotes composition with a reparametrization $\varphi$. If $\br$ has $k$ continuous derivatives we write $\br\in C^k$.}
\APS{局部研究曲线要用 Taylor 展开。若 $\br\in C^3$，则
\[
\br(t_0+h)=\br(t_0)+h\br'(t_0)+\frac{h^2}{2}\br''(t_0)+\frac{h^3}{6}\br'''(t_0)+o(h^3),
\]
每一阶导数都有几何含义：一阶给出切线方向，二阶给出弯曲方向和曲率，三阶给出挠率。带 Lagrange 余项的版本给出误差估计，这在第 2 章证明曲线论基本定理和一点的标准展开时都会用到。}{Local study of curves uses Taylor expansion. If $\br\in C^3$, then the displayed expansion holds, and each derivative has geometric meaning: the first gives the tangent direction, the second the direction of bending and the curvature, the third the torsion. The version with Lagrange remainder gives error estimates, and both the fundamental theorem of curve theory and the standard expansion at a point use it.}
\APS{向量函数的积分还给出弧长的来源：曲线从 $t_0$ 到 $t$ 的弧长是速率 $|\br'(u)|$ 的积分。因此弧长参数 $s$ 满足 $\dd s/\dd t=|\br'(t)|>0$，可以反解出 $t=t(s)$，得到单位速度参数化。这一步把"任意参数"换成"几何参数"，是后面所有曲率公式简洁的原因。}{Integration of vector functions also gives the origin of arc length: the arc length from $t_0$ to $t$ is the integral of the speed $|\br'(u)|$. Hence the arc-length parameter $s$ satisfies $\dd s/\dd t=|\br'(t)|>0$, can be inverted as $t=t(s)$, and yields a unit-speed parametrization. This step replaces an arbitrary parameter by a geometric one and is the reason all later curvature formulas are simple.}
\EX{取 $\br(t)=(\cos t,\sin t,t)$，则 $\br'=(-\sin t,\cos t,1)$、$|\br'|=\sqrt2$、$\br''=-(\cos t,\sin t,0)$、$\br'''=(\sin t,-\cos t,0)$。在 $t=0$ 处 $\br(0)=(1,0,0)$、$\br'(0)=(0,1,1)$、$\br''(0)=(-1,0,0)$、$\br'''(0)=(0,-1,0)$，于是
\[
\br(h)=(1,0,0)+h(0,1,1)+\frac{h^2}{2}(-1,0,0)+\frac{h^3}{6}(0,-1,0)+o(h^3).
\]
这正是一条螺旋线在一点附近的局部形状。}{Take $\br(t)=(\cos t,\sin t,t)$; then $\br'=(-\sin t,\cos t,1)$, $|\br'|=\sqrt2$, $\br''=-(\cos t,\sin t,0)$ and $\br'''=(\sin t,-\cos t,0)$. At $t=0$, $\br(0)=(1,0,0)$, $\br'(0)=(0,1,1)$, $\br''(0)=(-1,0,0)$ and $\br'''(0)=(0,-1,0)$, so the displayed expansion holds. This is precisely the local shape of a helix near one point.}
\APS{为什么曲线曲面论需要标架与向量函数？因为几何量必须用可微对象表达才能做分析。切向量、法向量、曲率都是标架导数的分量；曲线论基本定理和曲面论基本定理本质上是同一类问题：给定这些分量（$\kappa,\tau$ 或 $E,F,G,L,M,N$），是否存在相应的几何对象，且唯一到刚体运动。这背后是线性 ODE 组和可积性（Frobenius）理论，而理论的对象正是向量函数。}{Why does curve and surface theory need frames and vector functions? Because geometric quantities must be expressed through differentiable objects before analysis can begin. Tangent vectors, normal vectors and curvatures are components of derivatives of frames; the fundamental theorem of curve theory and that of surface theory are essentially the same type of question: given these components ($\kappa,\tau$ or $E,F,G,L,M,N$), does a geometric object exist, and is it unique up to rigid motion? Behind this lie linear ODE systems and integrability (Frobenius) theory, whose objects are exactly vector functions.}
\ASK{为什么用弧长参数而不是任意参数？}{Why use arc length rather than an arbitrary parameter?}{弧长参数使速度恒为 $1$，于是 $|T'|$ 直接就是曲率，Frenet 公式没有额外的 $|\bgamma'|$ 因子，曲率、挠率成为曲线在刚体运动下的不变量。任意参数化下公式要乘上 $|\bgamma'|$ 的幂次，计算和理论都变复杂；而且弧长是内蕴的几何量。}{The arc-length parameter has unit speed, so $|T'|$ is exactly the curvature and the Frenet formulas carry no extra factor $|\bgamma'|$; curvature and torsion become invariants under rigid motions. With an arbitrary parameter the formulas acquire powers of $|\bgamma'|$, complicating both computation and theory, and arc length itself is an intrinsic geometric quantity.}
\ASK{活动标架法与固定坐标系有什么区别？}{What is the difference between the moving-frame method and a fixed coordinate system?}{固定坐标系的基向量不变，但曲线在空间中的位置和方向都随之变化；活动标架的基向量附着在几何对象上，导数的分量直接给出曲率、挠率等几何信息。代价是基向量本身在变，必须同时记录标架的转动，这正是反对称矩阵的作用。}{In a fixed coordinate system the basis vectors are constant while the position and orientation of the curve change; in a moving frame the basis is attached to the geometric object, so the components of derivatives directly give geometric information such as curvature and torsion. The price is that the basis itself changes and its rotation must be recorded, which is exactly the role of the skew-symmetric matrix.}

\section[曲线论]{曲线论：Frenet 标架与基本定理 | Curve theory: Frenet frames and the fundamental theorem}
\subsection{正则参数曲线与弧长 | Regular curves and arc length}
\APS{空间曲线是区间 $I\subset\RR$ 到 $\RR^3$ 的可微映射 $\bgamma:I\to\RR^3$。若 $\bgamma\in C^k$（$k\ge1$）且 $\bgamma'(t)\ne0$ 对一切 $t$ 成立，则称 $\bgamma$ 是正则参数曲线；$\bgamma'(t)$ 给出切方向，$|\bgamma'(t)|$ 是速率。非正则点（如尖点）处切方向可能不存在或突变，曲线论的基本公式在那里失效，所以正则性是整套理论的前提。若 $\bgamma$ 只在一个小区间上正则，称局部正则曲线。}{A space curve is a differentiable map $\bgamma:I\to\RR^3$ from an interval $I\subset\RR$. If $\bgamma\in C^k$ with $k\ge1$ and $\bgamma'(t)\ne0$ for all $t$, it is called a regular parametrized curve; $\bgamma'(t)$ gives the tangent direction and $|\bgamma'(t)|$ the speed. At nonregular points, such as cusps, the tangent direction may fail to exist or may jump, and the basic formulas of curve theory break down, so regularity is a standing hypothesis. If $\bgamma$ is regular on a small interval only, it is called locally regular.}
\APS{同一几何曲线可以有无数参数化：若 $t=t(\tau)$ 是 $C^1$ 且 $\dd t/\dd\tau\ne0$，则 $\tilde\bgamma(\tau)=\bgamma(t(\tau))$ 与 $\bgamma$ 是同一曲线。弧长是最自然的参数：从固定点 $t_0$ 量起，
\[
s(t)=\int_{t_0}^t|\bgamma'(u)|\,\dd u,\qquad \frac{\dd s}{\dd t}=|\bgamma'(t)|>0 ,
\]
因此 $s$ 是 $t$ 的严格增函数，可以反解出 $t=t(s)$，得到单位速度参数化 $\bgamma(s)$，满足 $|\bgamma'(s)|=1$。}{The same geometric curve admits infinitely many parametrizations: if $t=t(\tau)$ is $C^1$ with $\dd t/\dd\tau\ne0$, then $\tilde\bgamma(\tau)=\bgamma(t(\tau))$ describes the same curve. Arc length is the most natural parameter: measured from a fixed point $t_0$, the displayed integral defines $s(t)$, and $\dd s/\dd t=|\bgamma'(t)|>0$. Hence $s$ is strictly increasing in $t$ and can be inverted as $t=t(s)$, producing the unit-speed parametrization $\bgamma(s)$ with $|\bgamma'(s)|=1$.}
\APS{弧长参数化把一个分析问题变成几何问题：曲线的切向量是单位向量，所有公式中不再出现速率因子。例如圆的弧长参数化是 $\bgamma(s)=r(\cos(s/r),\sin(s/r))$，满足 $|\bgamma'|=1$；而用角度 $t$ 参数化时 $|\bgamma'|=r$，曲率公式要除以速率。讲义后面默认曲线已用弧长参数，除非特别说明。}{Arc-length parametrization turns an analytic problem into a geometric one: the tangent vector is a unit vector and no speed factor appears in the formulas. For example the unit-speed parametrization of a circle is $\bgamma(s)=r(\cos(s/r),\sin(s/r))$, while parametrizing by the angle $t$ gives speed $r$ and requires dividing by the speed in curvature formulas. In the rest of these notes a curve is assumed to be arc-length parametrized unless stated otherwise.}

\subsection{曲率、挠率与 Frenet 公式 | Curvature, torsion and the Frenet formulas}
\APS{设 $\bgamma(s)$ 是单位速度曲线。切向量 $\bT(s)=\bgamma'(s)$ 是单位向量，其导数 $\bT'(s)$ 与 $\bT$ 正交，长度 $|\bT'|$ 度量切线方向随弧长转动的快慢，定义为曲率 $\kappa(s)=|\bT'(s)|\ge0$。当 $\kappa>0$ 时定义主法向量 $\bN=\bT'/\kappa$，再令副法向量 $\bB=\bT\times\bN$。三个单位向量两两正交且 $\bT\times\bN=\bB$，构成右手正交标架 $\{\bT,\bN,\bB\}$，称为沿曲线的 Frenet 标架。}{Let $\bgamma(s)$ be a unit-speed curve. The tangent $\bT(s)=\bgamma'(s)$ is a unit vector, so its derivative $\bT'(s)$ is orthogonal to $\bT$; the length $|\bT'|$ measures how fast the tangent direction turns with arc length and is defined to be the curvature $\kappa(s)=|\bT'(s)|\ge0$. When $\kappa>0$, define the principal normal $\bN=\bT'/\kappa$ and then the binormal $\bB=\bT\times\bN$. The three unit vectors are mutually orthogonal with $\bT\times\bN=\bB$, forming a right-handed orthonormal frame $\{\bT,\bN,\bB\}$ called the Frenet frame along the curve.}
\APS{曲率描述曲线偏离直线的程度：直线 $\kappa\equiv0$；圆 $\kappa=1/r$ 为常数。挠率描述曲线偏离平面的程度，定义为 $\tau=-\bB'\cdot\bN$（也可写 $\tau=\bB'\cdot\bN$，符号约定因书而异，本文采用前者）。$\tau=0$ 当且仅当 $\bB$ 为常向量，即曲线落在与 $\bB$ 垂直的平面内，因此平面曲线的挠率为零；挠率非零说明曲线不断扭出密切平面。}{Curvature measures deviation from a straight line: a line has $\kappa\equiv0$, and a circle has constant $\kappa=1/r$. Torsion measures deviation from a plane and is defined as $\tau=-\bB'\cdot\bN$ (some texts use the opposite sign; we adopt this one). One has $\tau=0$ exactly when $\bB$ is constant, that is when the curve lies in a plane perpendicular to $\bB$; hence plane curves have zero torsion, while nonzero torsion means the curve keeps twisting out of its osculating plane.}
\APS{对 $\bT,\bN,\bB$ 求导并用它们自身展开，就得到曲线论中最重要的公式。由于 $\bT'=\kappa\bN$ 已知，设 $\bN'=a\bT+b\bN+c\bB$；由 $\bN$ 是单位向量得 $b=0$，由 $\bN\cdot\bT=0$ 求导得 $a=-\kappa$，由定义 $\tau=-\bB'\cdot\bN$ 与 $\bB'=\bT'\times\bN+\bT\times\bN'$ 得 $c=\tau$。同理 $\bB'=-\tau\bN$。}{Differentiating $\bT,\bN,\bB$ and expanding in the frame itself yields the most important formulas of curve theory. Since $\bT'=\kappa\bN$ is known, write $\bN'=a\bT+b\bN+c\bB$; $b=0$ because $\bN$ is a unit vector, $a=-\kappa$ by differentiating $\bN\cdot\bT=0$, and $c=\tau$ from the definition $\tau=-\bB'\cdot\bN$ together with $\bB'=\bT'\times\bN+\bT\times\bN'$. Likewise $\bB'=-\tau\bN$.}
\FORMULA{Frenet 公式}{the Frenet formulas}
\[
\bT'=\kappa\bN,\qquad \bN'=-\kappa\bT+\tau\bB,\qquad \bB'=-\tau\bN .
\]
\READ{``T prime equals kappa N; N prime equals minus kappa T plus tau B; B prime equals minus tau N.''}
\APS{Frenet 公式说明 Frenet 标架的导数矩阵是反对称的：
\[
\begin{pmatrix}\bT\\ \bN\\ \bB\end{pmatrix}'
=\begin{pmatrix}0&\kappa&0\\ -\kappa&0&\tau\\ 0&-\tau&0\end{pmatrix}
\begin{pmatrix}\bT\\ \bN\\ \bB\end{pmatrix}.
\]
这正是第 1 章"标架导数是反对称矩阵乘以标架"的具体化：曲率和挠率就是标架转动的两个独立角速度分量。}{The Frenet formulas say that the derivative matrix of the Frenet frame is skew-symmetric, as displayed. This is exactly the concrete form of the Chapter 1 principle that the derivative of a frame is a skew-symmetric matrix times the frame: curvature and torsion are the two independent angular velocities of the frame's rotation.}
\APS{当曲线用任意参数 $t$ 给出时，曲率和挠率可用公式直接计算，不必先求弧长参数：
\[
\kappa=\frac{|\bgamma'\times\bgamma''|}{|\bgamma'|^3},\qquad
\tau=\frac{\det(\bgamma',\bgamma'',\bgamma''')}{|\bgamma'\times\bgamma''|^2}.
\]
第一式来自 $\bT'=(\bgamma''|\bgamma'|^2-\bgamma'(\bgamma'\cdot\bgamma''))/|\bgamma'|^4$ 与 $\bT'=\kappa\bN$；第二式由 $\bB=(\bgamma'\times\bgamma'')/|\bgamma'\times\bgamma''|$ 求导后与 $\bN$ 作内积得到。}{When the curve is given with an arbitrary parameter $t$, curvature and torsion can be computed directly without first finding arc length, by the displayed formulas. The first follows from the stated expression for $\bT'$ together with $\bT'=\kappa\bN$; the second comes from differentiating $\bB=(\bgamma'\times\bgamma'')/|\bgamma'\times\bgamma''|$ and taking the inner product with $\bN$.}
\EX{螺旋线 $\bgamma(t)=(a\cos t,a\sin t,bt)$（$a>0$）的完整计算。先求
\[
\bgamma'=(-a\sin t,a\cos t,b),\quad |\bgamma'|=\sqrt{a^2+b^2}=:c,\quad s=ct,
\]
故 $\bT=\frac1c(-a\sin t,a\cos t,b)$；再求 $\bgamma''=(-a\cos t,-a\sin t,0)$，
\[
\bT'=\frac1c\bgamma''=-\frac{a}{c}(\cos t,\sin t,0),\qquad
\kappa=|\bT'|=\frac{a}{c^2}=\frac{a}{a^2+b^2},
\]
于是 $\bN=-(\cos t,\sin t,0)$。计算外积得
\[
\bB=\bT\times\bN=\frac1c(b\sin t,-b\cos t,a),\qquad
\bB'=\frac1c(b\cos t,b\sin t,0)=-\frac{b}{c}\bN,
\]
故 $\tau=-\bB'\cdot\bN=b/c^2=b/(a^2+b^2)$。两个结果关于 $a,b$ 对称：$a=b=1$ 时 $\kappa=\tau=1/2$；$b=0$ 时 $\tau=0$，曲线退化为平面圆且 $\kappa=1/a$；$a\to0$ 时曲线拉直成 $z$ 轴，$\kappa,\tau\to0$。}{Full computation for the helix $\bgamma(t)=(a\cos t,a\sin t,bt)$ with $a>0$. First the displayed velocity and speed $c=\sqrt{a^2+b^2}$ give $s=ct$ and $\bT=(-a\sin t,a\cos t,b)/c$. Since $\bgamma''=(-a\cos t,-a\sin t,0)$, one gets $\bT'=-\frac{a}{c}(\cos t,\sin t,0)$ and hence $\kappa=|\bT'|=a/c^2=a/(a^2+b^2)$, with $\bN=-(\cos t,\sin t,0)$. The cross product gives the displayed $\bB$ and $\bB'=-(b/c)\bN$, hence $\tau=-\bB'\cdot\bN=b/c^2=b/(a^2+b^2)$. The two results are symmetric in $a,b$: for $a=b=1$, $\kappa=\tau=1/2$; for $b=0$, $\tau=0$ and the curve degenerates to a plane circle with $\kappa=1/a$; as $a\to0$ the curve straightens to the $z$-axis and both invariants tend to zero.}
\APS{这个例子说明了一条重要原则：两个不变量刻画曲线的局部形状。螺旋线在刚体运动下由 $\kappa,\tau$ 唯一决定；改变 $a,b$ 的比例只改变"弯曲与扭转的比"。在面试中若被要求算具体例子，按"速度—单位切向—曲率—主法向—副法向—挠率"的顺序写，分数最稳。}{This example illustrates an important principle: two invariants characterize the local shape of a curve. A helix is determined up to rigid motion by $\kappa,\tau$, and changing the ratio of $a$ to $b$ only changes the ratio of bending to twisting. In an interview, when asked for a concrete computation, follow the order speed, unit tangent, curvature, principal normal, binormal, torsion; this is the safest way to earn full marks.}

\subsection{曲线论基本定理 | The fundamental theorem of curve theory}
\APS{曲线论基本定理包含存在性和唯一性两部分。唯一性：设两条单位速度曲线 $\bgamma_1,\bgamma_2$ 定义在同一区间上，曲率都为正且对应点的 $\kappa,\tau$ 相同，则存在刚体运动 $g(\bx)=A\bx+\bb$（$A$ 为三阶正交矩阵、$\det A=1$）使 $\bgamma_2(s)=g(\bgamma_1(s))$。存在性：给定区间上连续函数 $\kappa>0$ 与 $\tau$，存在单位速度曲线以它们为曲率和挠率，且这样的曲线在刚体运动下唯一。}{The fundamental theorem of curve theory has a uniqueness and an existence part. Uniqueness: if two unit-speed curves $\bgamma_1,\bgamma_2$ are defined on the same interval, have positive curvature, and have the same $\kappa,\tau$ at corresponding points, then there is a rigid motion $g(\bx)=A\bx+\bb$ with $A$ orthogonal and $\det A=1$ such that $\bgamma_2(s)=g(\bgamma_1(s))$. Existence: given continuous functions $\kappa>0$ and $\tau$ on an interval, there is a unit-speed curve having them as curvature and torsion, and such a curve is unique up to rigid motion.}
\APS{唯一性的证明思路是"比较两个 Frenet 标架"。把两条曲线的标架写成矩阵 $F_i=(\bT_i,\bN_i,\bB_i)$，Frenet 公式给出 $F_i'=F_i\Omega$，其中 $\Omega$ 是同一个反对称矩阵。令 $A(s)=F_2(s)F_1(s)^{T}$，则 $A'=F_2'F_1^{T}+F_2(F_1')^{T}=F_2(\Omega-\Omega)F_1^{T}=0$，故 $A$ 为常正交矩阵。再对 $\bgamma_2-A\bgamma_1$ 求导并用标架比较，得其为常向量，于是 $\bgamma_2=A\bgamma_1+\bb$。}{The proof of uniqueness compares the two Frenet frames. Write each frame as a matrix $F_i=(\bT_i,\bN_i,\bB_i)$; the Frenet formulas give $F_i'=F_i\Omega$ with the same skew-symmetric matrix $\Omega$. Setting $A(s)=F_2(s)F_1(s)^{T}$ gives $A'=F_2'F_1^{T}+F_2(F_1')^{T}=F_2(\Omega-\Omega)F_1^{T}=0$, so $A$ is a constant orthogonal matrix. Differentiating $\bgamma_2-A\bgamma_1$ and comparing frames shows that it is a constant vector, hence $\bgamma_2=A\bgamma_1+\bb$.}
\APS{存在性用线性 ODE 组。以任意初始标架 $F(0)=I$ 解 Frenet 方程 $F'=F\Omega(s)$；这是系数连续的线性方程组，由 Picard--Lindelöf 定理在整区间上有唯一解，且解保持正交性（因为 $\Omega$ 反对称）。再令 $\bgamma(s)=\int_0^s\bT(u)\dd u+\bgamma(0)$，则 $\bgamma'=\bT$ 是单位向量，且按构造 $\bT'=\kappa\bN$、$\bN'=-\kappa\bT+\tau\bB$、$\bB'=-\tau\bN$，故 $\kappa,\tau$ 正是所求。}{Existence uses a linear ODE system. Solve the Frenet equations $F'=F\Omega(s)$ with initial frame $F(0)=I$; this is a linear system with continuous coefficients, so Picard--Lindelöf gives a unique solution on the whole interval, and the solution stays orthogonal because $\Omega$ is skew-symmetric. Then set $\bgamma(s)=\int_0^s\bT(u)\dd u+\bgamma(0)$; by construction $\bgamma'=\bT$ is a unit vector and the Frenet equations for $\bT,\bN,\bB$ hold, so $\kappa,\tau$ are exactly the prescribed functions.}
\APS{定理的意义是：曲线在刚体运动下的局部形状完全由两个函数 $\kappa(s),\tau(s)$ 决定，这与"平面曲线由曲率函数决定"是同一思想在三维的推广。它也解释了为什么 Frenet 公式比任何具体坐标公式都重要：它是把几何对象还原出来的 ODE。}{The meaning of the theorem is that the local shape of a curve up to rigid motion is completely determined by the two functions $\kappa(s),\tau(s)$, the three-dimensional analogue of the fact that a plane curve is determined by its curvature function. It also explains why the Frenet formulas matter more than any explicit coordinate formula: they are the ODE that reconstructs the geometric object.}

\subsection{一点的标准展开与曲线偶 | Standard expansion and pairs of curves}
\APS{把 $\bgamma(s)$ 在 $s=0$ 处 Taylor 展开，并用 Frenet 公式把各阶导数换成 $\bT_0,\bN_0,\bB_0$，得到一点的标准展开
\[
\bgamma(s)=\bgamma(0)+\left(s-\frac{\kappa^2s^3}{6}\right)\bT_0
+\left(\frac{\kappa s^2}{2}+\frac{\kappa's^3}{6}\right)\bN_0
+\frac{\kappa\tau s^3}{6}\bB_0+o(s^3),
\]
其中 $\kappa,\tau,\kappa'$ 均在 $s=0$ 处取值。}{Expand $\bgamma(s)$ in Taylor series at $s=0$ and use the Frenet formulas to express the derivatives in $\bT_0,\bN_0,\bB_0$; this gives the displayed standard expansion, where $\kappa,\tau,\kappa'$ are evaluated at $s=0$.}
\APS{这个展开有清晰的几何读法：沿 $\bT_0$ 的一阶项就是切线；沿 $\bN_0$ 的二阶项 $\frac{\kappa s^2}{2}$ 说明曲线向主法向弯曲，弯曲的强弱由 $\kappa$ 给出，这就是密切圆半径 $1/\kappa$；沿 $\bB_0$ 的最低阶项是三次的 $\frac{\kappa\tau s^3}{6}$，说明曲线离开密切平面的速度由挠率控制。展开到三阶，曲线在密切平面上的投影近似为抛物线，在法平面上的投影近似为三次曲线。}{The expansion has a clear geometric reading: the first-order term along $\bT_0$ is the tangent; the second-order term $\frac{\kappa s^2}{2}$ along $\bN_0$ says that the curve bends toward the principal normal with strength $\kappa$, which is the curvature of the osculating circle of radius $1/\kappa$; the lowest-order term along $\bB_0$ is the cubic $\frac{\kappa\tau s^3}{6}$, saying that the speed at which the curve leaves the osculating plane is controlled by the torsion. To third order the projection onto the osculating plane is approximately a parabola, and the projection onto the normal plane is approximately a cubic.}
\APS{第 2.7 节讨论"存在对应关系的曲线偶"，最典型的是 Bertrand 曲线：两条曲线在对应点有相同的主法线。可以证明一对正则曲线互为 Bertrand 对应，当且仅当存在常数 $a,b$（$a\ne0$）使 $a\kappa+b\tau=1$；圆柱螺旋线是 Bertrand 曲线，与它同轴、同螺距、相位相差 $\pi$ 的螺旋线构成对应偶。这个结论说明 $\kappa,\tau$ 不仅决定单条曲线，也决定曲线之间的特殊对应。}{Section 2.7 discusses pairs of curves in a correspondence; the standard example is a pair of Bertrand curves, two curves whose principal normals coincide at corresponding points. One can prove that two regular curves are in a Bertrand correspondence exactly when there are constants $a,b$ with $a\ne0$ and $a\kappa+b\tau=1$; a circular helix is a Bertrand curve, paired with the coaxial helix of the same pitch shifted in phase by $\pi$. This shows that $\kappa,\tau$ govern not only a single curve but also special correspondences between curves.}

\subsection{平面曲线 | Plane curves}
\APS{平面曲线是挠率为零的特例。设 $\bgamma(s)=(x(s),y(s))$ 为单位速度，则 $\bT=(x',y')=(\cos\theta,\sin\theta)$，其中 $\theta(s)$ 是切线与 $x$ 轴的夹角。定义相对曲率（带符号曲率）
\[
\kappa_r=\frac{\dd\theta}{\dd s}=\frac{x'y''-y'x''}{(x'^2+y'^2)^{3/2}},
\]
则 $\kappa=|\kappa_r|$，且 Frenet 公式变为 $\bT'=\kappa_r\bN$、$\bN'=-\kappa_r\bT$，其中 $\bN=(-\sin\theta,\cos\theta)$ 是 $\bT$ 逆时针转 $90^\circ$ 得到的法向量。$\kappa_r>0$ 表示曲线向左（逆时针）弯，$\kappa_r<0$ 表示向右弯；直线 $\kappa_r=0$，圆 $\kappa_r=\pm1/r$。}{A plane curve is the special case of zero torsion. Let $\bgamma(s)=(x(s),y(s))$ be unit speed; then $\bT=(x',y')=(\cos\theta,\sin\theta)$, where $\theta(s)$ is the angle between the tangent and the $x$-axis. Define the relative (signed) curvature by the displayed formula; then $\kappa=|\kappa_r|$ and the Frenet formulas become $\bT'=\kappa_r\bN$, $\bN'=-\kappa_r\bT$ with $\bN=(-\sin\theta,\cos\theta)$ obtained by rotating $\bT$ counterclockwise by $90^\circ$. Positive $\kappa_r$ means the curve turns left (counterclockwise) and negative means it turns right; a line has $\kappa_r=0$ and a circle has $\kappa_r=\pm1/r$.}
\APS{平面曲线的基本定理同样成立：给定连续函数 $\kappa_r(s)$，存在平面曲线以它为相对曲率，且唯一到平面刚体运动（平移与旋转）。证明只需对 $\theta(s)=\int\kappa_r\dd s$ 积分再积分一次得到曲线；这与三维情形相比少了挠率，因此平面曲线的形状只由一个函数决定。}{The fundamental theorem also holds for plane curves: given a continuous function $\kappa_r(s)$, there is a plane curve having it as relative curvature, unique up to planar rigid motions (translations and rotations). The proof only integrates $\theta(s)=\int\kappa_r\dd s$ and then integrates once more to get the curve; compared with the space case there is no torsion, so the shape of a plane curve is determined by a single function.}
\APS{曲率圆与渐屈线给出平面曲线的整体描述。曲率圆是与曲线在一点有二阶切触的圆，圆心为 $\bc(s)=\bgamma(s)+\bN(s)/\kappa_r(s)$，半径 $1/|\kappa_r|$；这些圆心的轨迹叫做渐屈线，它是曲线法线的包络。反过来，原曲线是渐屈线的一条渐伸线：渐伸线的切线是原曲线的法线，二者构成"法线族"的两种描述。若曲线是简单闭凸曲线，则相对曲率至少有四个极值点（四顶点定理）。}{The osculating circle and the evolute give a global description of a plane curve. The osculating circle is the circle having second-order contact with the curve at a point; its center is $\bc(s)=\bgamma(s)+\bN(s)/\kappa_r(s)$ and its radius is $1/|\kappa_r|$. The locus of these centers is the evolute, the envelope of the normals of the curve. Conversely, the original curve is an involute of the evolute: the tangent to an involute is a normal to the original curve, and the two describe the same family of normals in two ways. If the curve is a simple closed convex curve, its relative curvature has at least four extremal points, which is the four-vertex theorem.}
\EX{抛物线 $y=x^2$ 的相对曲率。取 $x$ 为参数，$x'=1$、$y'=2x$、$x''=0$、$y''=2$，故
\[
\kappa_r(x)=\frac{2}{(1+4x^2)^{3/2}},\qquad
\kappa_r(0)=2,\qquad \kappa_r(1)=\frac{2}{5\sqrt5}\approx0.1789 .
\]
在原点曲率半径为 $1/2$，曲率圆为 $(x)^2+(y-\frac12)^2=\frac14$；在 $x=1$ 处曲线几乎变直，曲率显著变小。这个算例也说明：曲率是局部量，随点变化。}{Relative curvature of the parabola $y=x^2$. Using $x$ as parameter, $x'=1$, $y'=2x$, $x''=0$, $y''=2$, hence the displayed relative curvature, with $\kappa_r(0)=2$ and $\kappa_r(1)=2/(5\sqrt5)\approx0.1789$. At the origin the radius of curvature is $1/2$ and the osculating circle is $x^2+(y-\tfrac12)^2=\tfrac14$; at $x=1$ the curve is nearly straight and the curvature has dropped considerably. The example also shows that curvature is a local quantity varying from point to point.}
\EX{圆与直线的对照。圆 $\bgamma(s)=r(\cos(s/r),\sin(s/r))$ 的切向量与曲率为
\[
\bT=(-\sin(s/r),\cos(s/r)),\qquad
\bT'=-\frac1r(\cos(s/r),\sin(s/r)),\qquad \kappa_r=-\frac1r,
\]
故 $|\kappa|=1/r$ 为常数（逆时针参数化时法向指向圆心，符号依定向而定）；直线 $\bgamma(s)=\ba s+\bb$ 的 $\bT=\ba$ 为常向量，故 $\kappa_r\equiv0$。这说明曲率确实"度量弯曲"，而相对曲率的符号记录弯曲方向。}{Circle versus line. For the circle $\bgamma(s)=r(\cos(s/r),\sin(s/r))$ the tangent and its derivative are the displayed vectors, so $\kappa_r=-1/r$ and $|\kappa|=1/r$ is constant (for the counterclockwise parametrization the normal points to the center, with the sign depending on orientation); for the line $\bgamma(s)=\ba s+\bb$, $\bT=\ba$ is constant and $\kappa_r\equiv0$. Curvature therefore does measure bending, while the sign of the relative curvature records the direction of bending.}
\ASK{曲率和挠率的几何意义是什么？}{What is the geometric meaning of curvature and torsion?}{曲率是切向量方向对弧长的转动率，$\kappa=|\bT'|$，衡量曲线偏离直线的程度，倒数 $1/\kappa$ 是曲率半径；挠率是副法向量（等价地密切平面）绕切线的转动率，$\tau=-\bB'\cdot\bN$，衡量曲线偏离平面的程度。$\kappa\equiv0$ 为直线，$\tau\equiv0$ 为平面曲线。}{Curvature is the rate at which the tangent direction turns with arc length, $\kappa=|\bT'|$; it measures deviation from a straight line, and $1/\kappa$ is the radius of curvature. Torsion is the rate at which the binormal, equivalently the osculating plane, rotates about the tangent, $\tau=-\bB'\cdot\bN$, and measures deviation from a plane. A line has $\kappa\equiv0$ and a plane curve has $\tau\equiv0$.}
\ASK{曲线论基本定理为什么重要？}{Why is the fundamental theorem of curve theory important?}{它说 $\kappa,\tau$ 是曲线在刚体运动下的完全不变量：给定两个连续函数就有曲线，且曲线只由它们决定。这把"研究曲线形状"化为"研究两个函数"，也是标架方法、ODE 存在唯一性在几何中的直接应用；曲面论基本定理是同一思想在二维的推广。}{It says that $\kappa,\tau$ are complete invariants of a curve up to rigid motion: given two continuous functions a curve exists and is determined by them. This reduces the study of curve shapes to the study of two functions and is a direct application of the moving-frame method and ODE existence-uniqueness in geometry; the fundamental theorem of surface theory is the two-dimensional analogue.}

\section[曲面的第一基本形式]{曲面的第一基本形式 | The first fundamental form of a surface}
\subsection{正则参数曲面与切平面 | Regular parametrized surfaces and tangent planes}
\APS{曲面的参数表示是定义在平面区域 $D$ 上的双参数向量函数 $\br(u,v)=(x,y,z)$，其中 $x,y,z$ 都是 $u,v$ 的函数。若 $\br\in C^k$（$k\ge1$）且偏导数向量 $\br_u,\br_v$ 线性无关，即 $\br_u\times\br_v\ne0$，则称 $\br$ 是正则参数曲面；$u,v$ 叫参数或坐标，固定 $u$ 或 $v$ 得到的曲线叫参数曲线（坐标曲线）。正则性保证每点有唯一切平面，曲面没有尖点或自交的"折痕"。}{A surface is a two-parameter vector function $\br(u,v)=(x,y,z)$ on a plane domain $D$, where $x,y,z$ are functions of $u,v$. If $\br\in C^k$ with $k\ge1$ and the partial derivative vectors $\br_u,\br_v$ are linearly independent, that is $\br_u\times\br_v\ne0$, then $\br$ is called a regular parametrized surface; $u,v$ are parameters or coordinates, and fixing $u$ or $v$ gives the parameter curves (coordinate curves). Regularity guarantees a unique tangent plane at every point and rules out cusps or folded edges.}
\APS{常见的曲面例子：平面 $\br(u,v)=(u,v,0)$；图像曲面 $\br(u,v)=(u,v,f(u,v))$，其中 $\br_u\times\br_v=(-f_u,-f_v,1)\ne0$；球面 $\br(u,v)=R(\sin u\cos v,\sin u\sin v,\cos u)$，在 $0<u<\pi$ 上正则；环面 $\br(u,v)=((R+r\cos u)\cos v,(R+r\cos u)\sin v,r\sin u)$，其中 $R>r>0$；旋转面 $\br(u,v)=(f(u)\cos v,f(u)\sin v,g(u))$，$f>0$。这些例子覆盖了课程中绝大多数计算。}{Common examples of surfaces: the plane $\br(u,v)=(u,v,0)$; a graph $\br(u,v)=(u,v,f(u,v))$ with $\br_u\times\br_v=(-f_u,-f_v,1)\ne0$; the sphere $\br(u,v)=R(\sin u\cos v,\sin u\sin v,\cos u)$, regular for $0<u<\pi$; the torus $\br(u,v)=((R+r\cos u)\cos v,(R+r\cos u)\sin v,r\sin u)$ with $R>r>0$; and a surface of revolution $\br(u,v)=(f(u)\cos v,f(u)\sin v,g(u))$ with $f>0$. These examples cover most computations in the course.}
\APS{在正则点 $p=\br(u,v)$ 处，切平面是 $\br_u,\br_v$ 张成的二维平面
\[
T_pS=\operatorname{span}\{\br_u,\br_v\}=\{\br_u a+\br_v b:\ a,b\in\RR\},
\]
单位法向量为 $\bn=(\br_u\times\br_v)/|\br_u\times\br_v|$。曲面上一条曲线 $\bgamma(t)=\br(u(t),v(t))$ 的切向量是 $\bgamma'=\br_u u'+\br_v v'$，必在切平面内，这是"切向量"名称的来源。若曲面由隐式方程 $F(x,y,z)=0$ 给出，则法向平行于 $\nabla F$，因为沿曲面 $F$ 的导数恒为零。}{At a regular point $p=\br(u,v)$ the tangent plane is the plane spanned by $\br_u,\br_v$, displayed, with unit normal $\bn=(\br_u\times\br_v)/|\br_u\times\br_v|$. The tangent vector of a curve $\bgamma(t)=\br(u(t),v(t))$ on the surface is $\bgamma'=\br_u u'+\br_v v'$, which necessarily lies in the tangent plane; this is where the name comes from. If the surface is given implicitly by $F(x,y,z)=0$, the normal is parallel to $\nabla F$, since the derivative of $F$ along the surface vanishes identically.}

\subsection{第一基本形式 | The first fundamental form}
\APS{把切向量的长度用参数增量表示，就得到第一基本形式。设 $\dd\br=\br_u\dd u+\br_v\dd v$，则
\[
I=\dd\br\cdot\dd\br=E\,\dd u^2+2F\,\dd u\,\dd v+G\,\dd v^2,
\]
其中
\[
E=\br_u\cdot\br_u,\qquad F=\br_u\cdot\br_v,\qquad G=\br_v\cdot\br_v .
\]
由 $\br_u,\br_v$ 线性无关得 $E>0$、$G>0$、$EG-F^2=|\br_u\times\br_v|^2>0$，所以 $I$ 是正定二次型，可以看作曲面上的"度量张量"。}{Expressing the length of a tangent vector through parameter increments gives the first fundamental form. With $\dd\br=\br_u\dd u+\br_v\dd v$, one obtains the displayed $I=E\dd u^2+2F\dd u\dd v+G\dd v^2$, where $E=\br_u\cdot\br_u$, $F=\br_u\cdot\br_v$ and $G=\br_v\cdot\br_v$. Linear independence of $\br_u,\br_v$ gives $E>0$, $G>0$ and $EG-F^2=|\br_u\times\br_v|^2>0$, so $I$ is a positive definite quadratic form, the metric tensor of the surface.}
\APS{第一基本形式给出曲面上的三种基本测量。弧长：曲线 $u=u(t),v=v(t)$（$a\le t\le b$）的弧长是 $\int_a^b\sqrt{Eu'^2+2Fu'v'+Gv'^2}\,\dd t$。夹角：两个切向量 $\dd\br_1,\dd\br_2$ 的夹角余弦为 $I(\dd\br_1,\dd\br_2)/(|\dd\br_1||\dd\br_2|)$；特别地两条坐标曲线的夹角 $\omega$ 满足 $\cos\omega=F/\sqrt{EG}$。面积：参数区域 $D$ 对应的曲面面积为
\[
A=\iint_D\sqrt{EG-F^2}\,\dd u\,\dd v .
\]
这三个公式是内蕴几何的出发点：它们只用 $E,F,G$，不涉及曲面在空间中的位置。}{The first fundamental form provides the three basic measurements on a surface. Arc length: the length of a curve $u=u(t),v=v(t)$ is the displayed integral. Angle: the cosine of the angle between tangent vectors $\dd\br_1,\dd\br_2$ is $I(\dd\br_1,\dd\br_2)/(|\dd\br_1||\dd\br_2|)$; in particular the angle $\omega$ between the two coordinate curves satisfies $\cos\omega=F/\sqrt{EG}$. Area: the area corresponding to a parameter domain $D$ is the displayed double integral. These three formulas are the starting point of intrinsic geometry: they use only $E,F,G$ and not the position of the surface in space.}
\APS{若 $F=0$，称参数曲线网是正交的。此时 $I=E\dd u^2+G\dd v^2$，弧长元为 $\sqrt E|\dd u|$ 与 $\sqrt G|\dd v|$ 的平方和，面积元为 $\sqrt{EG}\,\dd u\dd v$，计算大为简化。第 3.4 节证明：在曲面的非脐点附近总可以取到正交参数曲线网；更强地，局部总存在等温坐标，使 $E=G=\lambda^2$、$F=0$，即 $I=\lambda^2(\dd u^2+\dd v^2)$。等温坐标的存在性归结为解一个 Beltrami 方程，需要复分析或椭圆型 PDE 理论；它说明任意曲面局部都"保角于平面"。}{If $F=0$ the parameter net is called orthogonal. Then $I=E\dd u^2+G\dd v^2$, the arc element is the square root of the sum of $\sqrt E|\dd u|$ and $\sqrt G|\dd v|$ squared, and the area element is $\sqrt{EG}\,\dd u\dd v$, which greatly simplifies computations. Section 3.4 proves that near any non-umbilic point one can always choose an orthogonal parameter net; more strongly, isothermal coordinates always exist locally, with $E=G=\lambda^2$ and $F=0$, that is $I=\lambda^2(\dd u^2+\dd v^2)$. Existence of isothermal coordinates reduces to a Beltrami equation and needs complex analysis or elliptic PDE theory; it says that every surface is locally conformal to a plane.}

\subsection{第一基本形式的完整计算 | Complete computations of first fundamental forms}
\EX{球面 $\br(u,v)=R(\sin u\cos v,\sin u\sin v,\cos u)$。先求偏导数
\[
\br_u=R(\cos u\cos v,\cos u\sin v,-\sin u),\qquad
\br_v=R(-\sin u\sin v,\sin u\cos v,0).
\]
直接计算内积：$E=R^2(\cos^2u\cos^2v+\cos^2u\sin^2v+\sin^2u)=R^2$；$F=R^2(-\cos u\sin u\cos v\sin v+\cos u\sin u\sin v\cos v)=0$；$G=R^2(\sin^2u\sin^2v+\sin^2u\cos^2v)=R^2\sin^2u$。故
\[
I=R^2\dd u^2+R^2\sin^2u\,\dd v^2,
\]
面积 $A=\int_0^{2\pi}\int_0^\pi R^2\sin u\,\dd u\dd v=4\pi R^2$，与 Euclid 球面积公式一致。}{Sphere $\br(u,v)=R(\sin u\cos v,\sin u\sin v,\cos u)$. The displayed partial derivatives give, by direct inner products, $E=R^2$, $F=0$ and $G=R^2\sin^2u$, hence the displayed first fundamental form. The area is $A=\int_0^{2\pi}\int_0^\pi R^2\sin u\,\dd u\dd v=4\pi R^2$, consistent with the Euclidean formula for the area of a sphere.}
\EX{环面 $\br(u,v)=((R+r\cos u)\cos v,(R+r\cos u)\sin v,r\sin u)$，$R>r>0$。偏导数为
\[
\br_u=r(-\sin u\cos v,-\sin u\sin v,\cos u),\qquad
\br_v=(R+r\cos u)(-\sin v,\cos v,0),
\]
两者正交，且 $|\br_u|=r$、$|\br_v|=R+r\cos u$，故
\[
I=r^2\dd u^2+(R+r\cos u)^2\dd v^2 .
\]
面积 $A=\int_0^{2\pi}\int_0^{2\pi}r(R+r\cos u)\,\dd u\dd v=4\pi^2Rr$，这正是环面的经典面积公式。}{Torus $\br(u,v)=((R+r\cos u)\cos v,(R+r\cos u)\sin v,r\sin u)$ with $R>r>0$. The displayed partial derivatives are orthogonal with $|\br_u|=r$ and $|\br_v|=R+r\cos u$, hence the displayed first fundamental form. The area is $A=\int_0^{2\pi}\int_0^{2\pi}r(R+r\cos u)\,\dd u\dd v=4\pi^2Rr$, the classical formula for the area of a torus.}
\APS{旋转面 $\br(u,v)=(f(u)\cos v,f(u)\sin v,g(u))$（$f>0$）是一般公式的来源。计算得
\[
E=f'^2+g'^2,\qquad F=0,\qquad G=f^2,\qquad
I=(f'^2+g'^2)\dd u^2+f^2\dd v^2 .
\]
球面取 $f=R\sin u$、$g=R\cos u$ 即回到 $I=R^2(\dd u^2+\sin^2u\,\dd v^2)$；悬链面取 $f=a\cosh(u/a)$、$g=u$ 得 $I=\cosh^2(u/a)(\dd u^2+a^2\dd v^2)$；柱面取 $f=R$、$g=u$ 得 $I=\dd u^2+R^2\dd v^2$。这个统一公式是后面计算主曲率和测地线的基础。}{A surface of revolution $\br(u,v)=(f(u)\cos v,f(u)\sin v,g(u))$ with $f>0$ is the source of a general formula. A direct computation gives $E=f'^2+g'^2$, $F=0$ and $G=f^2$, hence the displayed first fundamental form. The sphere is the case $f=R\sin u$, $g=R\cos u$, recovering $I=R^2(\dd u^2+\sin^2u\,\dd v^2)$; the catenoid is the case $f=a\cosh(u/a)$, $g=u$, giving $I=\cosh^2(u/a)(\dd u^2+a^2\dd v^2)$; the cylinder is the case $f=R$, $g=u$, giving $I=\dd u^2+R^2\dd v^2$. This unified formula is the basis for computing principal curvatures and geodesics later.}

\subsection{保长对应与保角对应 | Isometries and conformal maps}
\APS{两个曲面之间的映射若保持第一基本形式，称为保长对应（局部等距）。精确地说，若对应点处 $\dd\br_1\cdot\dd\br_1=\dd\tilde\br_1\cdot\dd\tilde\br_1$ 对一切切向量成立，则弧长、夹角、面积都被保持。最典型的例子是平面与柱面：单位柱面 $\br(u,v)=(\cos u,\sin u,v)$ 与平面 $\tilde\br(u,v)=(u,v,0)$ 的 $I$ 都是 $\dd u^2+\dd v^2$，把柱面沿母线剪开摊平就是等距映射。}{A map between two surfaces that preserves the first fundamental form is called an isometry (local isometry). Precisely, if $\dd\br_1\cdot\dd\br_1=\dd\tilde\br_1\cdot\dd\tilde\br_1$ at corresponding points for all tangent vectors, then arc length, angle and area are preserved. The standard example is the plane and the cylinder: the unit cylinder $\br(u,v)=(\cos u,\sin u,v)$ and the plane $\tilde\br(u,v)=(u,v,0)$ both have $I=\dd u^2+\dd v^2$, and cutting the cylinder along a generator and flattening it is an isometry.}
\APS{若映射把第一基本形式乘以一个正函数，即 $I'=\lambda^2I$，称为保角对应。保角映射保持夹角（因而保持形状），但长度和面积按 $\lambda$ 与 $\lambda^2$ 缩放。球极投影把球面（去掉北极）保角地映到平面；地图中的 Mercator 投影也是保角的：取球面 $I=R^2(\dd u^2+\sin^2u\,\dd v^2)$，令 $x=v$、$y=\int\dd u/\sin u=\ln\tan(u/2)$，则 $I=R^2\sin^2u(\dd x^2+\dd y^2)$，确实是平面的共形度量。}{If a map multiplies the first fundamental form by a positive function, $I'=\lambda^2I$, it is called conformal. A conformal map preserves angles, hence shapes, while lengths and areas scale by $\lambda$ and $\lambda^2$. Stereographic projection maps the sphere minus the north pole conformally to the plane; the Mercator projection used in cartography is also conformal: on the sphere with $I=R^2(\dd u^2+\sin^2u\,\dd v^2)$, set $x=v$ and $y=\int\dd u/\sin u=\ln\tan(u/2)$; then $I=R^2\sin^2u(\dd x^2+\dd y^2)$, a conformal metric of the plane.}
\APS{区分保长与保角对面试很重要：柱面与平面既保长又保角；球面与平面只能保角（球极投影、Mercator），不能保长，因为第 5 章的绝妙定理将证明 Gauss 曲率是保长不变量，而球面 $K=1/R^2$、平面 $K=0$。这个不可能性是最能体现"曲率是内蕴量"的例子。}{Distinguishing isometry from conformality matters in an interview: the cylinder and the plane are both isometric and conformal; the sphere and the plane are only conformal (stereographic projection, Mercator), never isometric, because the Theorema Egregium in Chapter 5 shows that the Gaussian curvature is an isometric invariant, while the sphere has $K=1/R^2$ and the plane has $K=0$. This impossibility is the best illustration that curvature is intrinsic.}

\subsection{可展曲面 | Developable surfaces}
\APS{直纹面是由一族直线（直母线）生成的曲面，参数形式为 $\br(u,v)=\ba(u)+v\,\bb(u)$，其中 $\ba(u)$ 是准线、$\bb(u)$ 是直母线的方向。可展曲面是沿每条直母线切平面都相同的直纹面，等价条件是 $\det(\ba',\bb,\bb')=0$。三类基本的可展曲面是柱面（$\bb$ 为常向量）、锥面（直母线过定点）和切线面（直母线是某条空间曲线的切线）。}{A ruled surface is generated by a one-parameter family of straight lines (generators) and has the form $\br(u,v)=\ba(u)+v\,\bb(u)$, where $\ba(u)$ is a directrix and $\bb(u)$ the direction of the generator. A developable surface is a ruled surface whose tangent plane is the same along each generator, equivalently $\det(\ba',\bb,\bb')=0$. The three basic types are cylinders ($\bb$ constant), cones (all generators pass through a fixed point) and tangent surfaces (the generators are tangents to a space curve).}
\APS{可展曲面的核心性质是 Gauss 曲率为零，且可以等距地展开成平面：沿直母线剪开后摊平，长度和角度都不改变。这正是制图、钣金和服装裁剪中"把曲面展开成平面"的数学依据。柱面与平面等距；锥面去掉顶点后与平面等距；而球面不能展开，因为 $K=1/R^2\ne0$。第 4、5 章会从第二基本形式和绝妙定理两个角度重新解释这件事。}{The central property of a developable surface is that its Gaussian curvature is zero and that it can be developed isometrically onto a plane: after cutting along a generator and flattening, lengths and angles are unchanged. This is the mathematical basis of unrolling a surface in cartography, sheet-metal work and garment cutting. Cylinders are isometric to the plane; cones with the vertex removed are isometric to the plane; the sphere cannot be developed because $K=1/R^2\ne0$. Chapters 4 and 5 explain this again from the viewpoints of the second fundamental form and the Theorema Egregium.}
\ASK{为什么第一基本形式是内蕴的？}{Why is the first fundamental form intrinsic?}{$E,F,G$ 只由曲面上的切向量内积决定，不依赖曲面如何嵌入 $\RR^3$；它们决定曲面上的长度、角度、面积，也决定测地线、平行移动和 Gauss 曲率。两个曲面即使外蕴形状不同（如平面与柱面），只要第一基本形式相同，内蕴几何就完全相同。}{$E,F,G$ are determined solely by inner products of tangent vectors on the surface and do not depend on how the surface is embedded in $\RR^3$; they determine length, angle and area on the surface, and also geodesics, parallel transport and Gaussian curvature. Two surfaces with different extrinsic shapes, such as a plane and a cylinder, have exactly the same intrinsic geometry whenever their first fundamental forms agree.}
\ASK{可展曲面为什么重要？}{Why are developable surfaces important?}{它们是把曲面局部等距展开成平面的对象，等价于 Gauss 曲率为零，也等价于沿直母线切平面不变。工程上对应可展开的板材、钣金件和地图；理论上它们是把内蕴几何与外蕴几何联系起来的最简单例子，并且是第 5 章绝妙定理的直接应用。}{They are the surfaces that can be developed isometrically onto a plane, equivalently they have zero Gaussian curvature and constant tangent plane along each generator. In engineering they correspond to unfoldable sheet metal and maps; theoretically they are the simplest link between intrinsic and extrinsic geometry and a direct application of the Theorema Egregium in Chapter 5.}

\section[第二基本形式]{曲面的第二基本形式与曲率 | The second fundamental form and curvature}
\subsection{第二基本形式 | The second fundamental form}
\APS{第一基本形式只记录曲面上的长度、角度和面积，无法区分"平"与"弯"：平面和球面在很小的局部都可以用相似的坐标描述，但它们的弯曲程度完全不同。描述弯曲需要看法向量如何变化，等价地看曲面上的曲线加速度的法向分量。设 $\bn$ 是单位法向量，定义
\[
L=\br_{uu}\cdot\bn,\qquad M=\br_{uv}\cdot\bn,\qquad N=\br_{vv}\cdot\bn,
\]
则第二基本形式为
\[
II=L\,\dd u^2+2M\,\dd u\,\dd v+N\,\dd v^2 .
\]
把法向量换成 $-\bn$，$II$ 变号，所以 $II$ 的符号依赖曲面的定向，但其零点、主曲率等几何信息不依赖。}{The first fundamental form records only length, angle and area on a surface and cannot distinguish flat from curved: a plane and a sphere look similar in a small neighbourhood with suitable coordinates, yet their bending is completely different. Describing bending requires looking at how the normal vector changes, equivalently at the normal component of the acceleration of curves on the surface. Let $\bn$ be the unit normal and define $L=\br_{uu}\cdot\bn$, $M=\br_{uv}\cdot\bn$, $N=\br_{vv}\cdot\bn$; then the second fundamental form is the displayed quadratic form. Replacing $\bn$ by $-\bn$ changes its sign, so the sign of $II$ depends on the orientation of the surface, while its zeros and the principal curvatures do not.}
\APS{第二基本形式有直接的几何解释。设 $\bgamma(t)=\br(u(t),v(t))$ 是曲面上一条曲线，则
\[
\bgamma''=\br_{uu}u'^2+2\br_{uv}u'v'+\br_{vv}v'^2+\br_u u''+\br_v v'',
\qquad
\bgamma''\cdot\bn=II(\bgamma',\bgamma').
\]
所以 $II$ 度量曲线加速度中垂直于曲面的部分：它与切向部分（由第一基本形式描述）一起，把曲面上的曲线运动分解为"沿曲面"和"离开曲面"两个方向。}{The second fundamental form has a direct geometric interpretation. For a curve $\bgamma(t)=\br(u(t),v(t))$ on the surface, the displayed decomposition holds, so $\bgamma''\cdot\bn=II(\bgamma',\bgamma')$. Thus $II$ measures the part of the acceleration perpendicular to the surface; together with the tangential part, described by the first fundamental form, it splits the motion of a curve on the surface into an along-surface and an off-surface component.}

\subsection{法曲率与 Weingarten 映射 | Normal curvature and the Weingarten map}
\APS{给定单位切向量 $\bv\in T_pS$，用过 $\bv$ 与法向量 $\bn$ 的平面去截曲面，得到一条法截线，其曲率称为沿 $\bv$ 的法曲率：
\[
\kappa_n(\bv)=\frac{II(\bv,\bv)}{I(\bv,\bv)} .
\]
法曲率只依赖切方向，不依赖具体曲线。Meusnier 定理说：曲面上过同一点、具有同一单位切向量的所有曲线，其曲率 $\kappa$ 与法曲率的关系是 $\kappa_n=\kappa\cos\varphi$，其中 $\varphi$ 是曲线主法向量与曲面法向量的夹角。}{Given a unit tangent vector $\bv\in T_pS$, intersect the surface with the plane through $\bv$ and the normal $\bn$; the resulting normal section has a curvature called the normal curvature in the direction $\bv$, given by the displayed quotient. The normal curvature depends only on the tangent direction, not on the particular curve. Meusnier's theorem says that for all curves on the surface through a point with the same unit tangent vector, the curvature $\kappa$ and the normal curvature satisfy $\kappa_n=\kappa\cos\varphi$, where $\varphi$ is the angle between the curve's principal normal and the surface normal.}
\APS{把法向量的变化写成切映射，就得到 Weingarten 映射
\[
S:T_pS\to T_pS,\qquad S(\bv)=-\dd\bn(\bv)=-\left.\frac{\dd}{\dd t}\bn(\bgamma(t))\right|_{t=0},
\]
其中 $\bgamma(0)=p$、$\bgamma'(0)=\bv$。$S$ 是自伴线性变换：$\langle S(\bv),\bw\rangle=\langle\bv,S(\bw)\rangle$，这等价于第二基本形式的对称性。因此在单位正交基下 $S$ 可以对角化，得到两个实特征值 $\kappa_1,\kappa_2$（主曲率）和两个正交特征方向（主方向）。}{Writing the change of the normal as a tangent map gives the Weingarten map displayed, where $\bgamma(0)=p$ and $\bgamma'(0)=\bv$. $S$ is a self-adjoint linear transformation, $\langle S(\bv),\bw\rangle=\langle\bv,S(\bw)\rangle$, which is equivalent to the symmetry of the second fundamental form. Hence in an orthonormal basis $S$ is diagonalizable, with two real eigenvalues $\kappa_1,\kappa_2$ (the principal curvatures) and two orthogonal eigendirections (the principal directions).}
\APS{在基 $\{\br_u,\br_v\}$ 下，$S$ 的矩阵是 $I^{-1}II$：
\[
(h^i_j)=\begin{pmatrix}E&F\\F&G\end{pmatrix}^{-1}
\begin{pmatrix}L&M\\M&N\end{pmatrix},
\qquad
II(\bv,\bw)=\langle S(\bv),\bw\rangle .
\]
于是 Gauss 曲率和平均曲率是 $S$ 的行列式与迹的一半：
\[
K=\det S=\frac{LN-M^2}{EG-F^2},\qquad
H=\frac12\operatorname{tr}S=\frac{EN-2FM+GL}{2(EG-F^2)} .
\]
这两个公式是曲面论中计算量最大的公式，必须背熟并能从 $I^{-1}II$ 推出。}{In the basis $\{\br_u,\br_v\}$ the matrix of $S$ is $I^{-1}II$, displayed, and $II(\bv,\bw)=\langle S(\bv),\bw\rangle$. Hence the Gaussian and mean curvatures are the determinant and half the trace of $S$, as displayed. These two formulas are the most heavily used computations in surface theory; they must be memorized and derivable from $I^{-1}II$.}
\APS{主曲率有明确的几何意义。若 $\kappa_1=\kappa_2$，称该点为脐点，此时所有方向的法曲率相同，球面上每一点都是脐点；若 $\kappa_1\ne\kappa_2$，两个主方向给出法曲率的极大与极小方向。Euler 公式给出任意方向的法曲率：若 $\bv$ 与第一主方向的夹角为 $\theta$，则
\[
\kappa_n(\theta)=\kappa_1\cos^2\theta+\kappa_2\sin^2\theta .
\]
沿主方向还有 Rodrigues 公式 $\dd\bn(\bv)=-\kappa\bv$，即法向量沿主方向的变化恰好是缩放。}{The principal curvatures have a clear geometric meaning. If $\kappa_1=\kappa_2$ the point is called umbilical, all directions have the same normal curvature, and every point of a sphere is umbilical; if $\kappa_1\ne\kappa_2$, the two principal directions give the directions of maximal and minimal normal curvature. Euler's formula gives the normal curvature in an arbitrary direction: if $\bv$ makes angle $\theta$ with the first principal direction, then the displayed formula holds. Along a principal direction Rodrigues' formula $\dd\bn(\bv)=-\kappa\bv$ says that the change of the normal is exactly a scaling.}

\subsection{主方向和主曲率的计算 | Computing principal directions and curvatures}
\APS{主曲率是特征方程 $\det(II-\kappa I)=0$ 的根：
\[
(EG-F^2)\kappa^2-(EN-2FM+GL)\kappa+(LN-M^2)=0 .
\]
由韦达定理，两根之和为 $2H$、之积为 $K$，这与上一节的公式一致。主方向满足 $(L-\kappa E)\dd u+(M-\kappa F)\dd v=0$；消去 $\kappa$ 得到主方向方程
\[
(EM-FL)\,\dd u^2+(EN-GL)\,\dd u\,\dd v+(FN-GM)\,\dd v^2=0 .
\]
当 $F=0$（正交参数网）时方程化为 $(EM-FL)\dd u^2+(EN-GL)\dd u\dd v-GM\dd v^2=0$，通常可分解为两个方向。}{The principal curvatures are the roots of the characteristic equation $\det(II-\kappa I)=0$, displayed. By Vieta's formulas their sum is $2H$ and their product is $K$, consistent with the previous subsection. Principal directions satisfy $(L-\kappa E)\dd u+(M-\kappa F)\dd v=0$; eliminating $\kappa$ gives the displayed principal-direction equation. When $F=0$ (orthogonal parameter net) the equation reduces to $(EM-FL)\dd u^2+(EN-GL)\dd u\dd v-GM\dd v^2=0$, which usually factors into two directions.}
\APS{一个实用的记忆方法是：主方向方程就是"把 $II-\kappa I$ 的行列式展开中的 $\kappa$ 消掉"，它关于 $\dd u,\dd v$ 是二次齐次式。若二次型退化（系数全为零），则该点为脐点，主方向不确定。计算时先算 $E,F,G,L,M,N$，再算 $K,H$，最后按需要求主方向；面试中若能当场从 $II-\kappa I$ 推出主方向方程，会显得理解透彻。}{A practical way to remember this is that the principal-direction equation is obtained by eliminating $\kappa$ from the determinant of $II-\kappa I$; it is a homogeneous quadratic in $\dd u,\dd v$. If the quadratic form degenerates (all coefficients vanish), the point is umbilical and the principal directions are undefined. In computations, first find $E,F,G,L,M,N$, then $K,H$, and finally the principal directions if needed; deriving the principal-direction equation on the spot from $II-\kappa I$ shows genuine understanding in an interview.}

\subsection{Dupin 标形与一点的标准展开 | The Dupin indicatrix and the standard expansion}
\APS{在一点 $p$ 取主方向为 $x,y$ 轴、法向量为 $z$ 轴，则曲面局部可写成 Monge 形式
\[
z=\frac12\left(\kappa_1x^2+\kappa_2y^2\right)+O(x^3,y^3),
\]
其中 $\kappa_1,\kappa_2$ 是 $p$ 处的主曲率。这个标准展开说明：二阶近似由两个主曲率完全决定，曲面在 $p$ 附近像一个"二阶曲面"。Dupin 标形是切平面内满足 $II(\bv,\bv)=\pm1$ 的点集，即
\[
\kappa_1x^2+\kappa_2y^2=\pm1 .
\]
}{At a point $p$ choose the principal directions as the $x,y$ axes and the normal as the $z$ axis; then the surface is locally the Monge form displayed, where $\kappa_1,\kappa_2$ are the principal curvatures at $p$. This standard expansion says that the second-order approximation is completely determined by the two principal curvatures, so the surface near $p$ behaves like a quadric. The Dupin indicatrix is the set of vectors in the tangent plane with $II(\bv,\bv)=\pm1$, that is $\kappa_1x^2+\kappa_2y^2=\pm1$.}
\APS{按 $K=\kappa_1\kappa_2$ 的符号可以把点分类：$K>0$ 时两个主曲率同号，Dupin 标形是椭圆，称为椭圆点，曲面局部位于切平面同侧（如球面）；$K<0$ 时异号，标形是双曲线，称为双曲点，曲面局部穿过切平面（如鞍面）；$K=0$ 时至少一个主曲率为零，标形退化为抛物线或一对平行线，称为抛物点（如柱面）。$K$ 的符号因此决定了曲面局部的凹凸类型，这是 Gauss 曲率最重要的定性含义。}{The sign of $K=\kappa_1\kappa_2$ classifies points: if $K>0$ the principal curvatures have the same sign, the Dupin indicatrix is an ellipse and the point is elliptic, with the surface locally on one side of its tangent plane (as on a sphere); if $K<0$ they have opposite signs, the indicatrix is a hyperbola and the point is hyperbolic, with the surface locally crossing its tangent plane (as on a saddle); if $K=0$ at least one principal curvature vanishes, the indicatrix degenerates to a parabola or a pair of parallel lines, and the point is parabolic (as on a cylinder). The sign of $K$ therefore determines the local convexity type of the surface, the most important qualitative meaning of the Gaussian curvature.}

\subsection{特殊曲面 | Special surfaces}
\APS{旋转面 $\br(u,v)=(f(u)\cos v,f(u)\sin v,g(u))$ 的两族主方向就是经线（$u$ 曲线）和纬线（$v$ 曲线）。由 $E=f'^2+g'^2$、$F=0$、$G=f^2$ 及第二基本形式可得两个主曲率
\[
\kappa_1=\frac{f'g''-f''g'}{(f'^2+g'^2)^{3/2}},\qquad
\kappa_2=\frac{g'}{f\sqrt{f'^2+g'^2}},
\]
从而 $K=\kappa_1\kappa_2$、$H=(\kappa_1+\kappa_2)/2$。经线方向弯曲由轮廓线的曲率决定，纬线方向的弯曲由旋转半径 $f$ 与轮廓线斜率 $g'$ 的比值决定。}{For a surface of revolution $\br(u,v)=(f(u)\cos v,f(u)\sin v,g(u))$ the two families of principal directions are the meridians ($u$-curves) and the parallels ($v$-curves). From $E=f'^2+g'^2$, $F=0$, $G=f^2$ and the second fundamental form one obtains the displayed principal curvatures, hence $K=\kappa_1\kappa_2$ and $H=(\kappa_1+\kappa_2)/2$. The bending in the meridian direction is governed by the curvature of the profile, and the bending in the parallel direction by the ratio of the rotation radius $f$ to the slope $g'$ of the profile.}
\APS{直纹面 $\br(u,v)=\ba(u)+v\,\bb(u)$ 的 Gauss 曲率可以写成
\[
K=-\frac{\det(\ba',\bb,\bb')^2}{EG-F^2}\le0 ,
\]
因此直纹面（除可展曲面外）是负曲率曲面；$K\equiv0$ 当且仅当 $\det(\ba',\bb,\bb')=0$，即曲面可展。双曲抛物面（马鞍面）与单叶双曲面是负曲率的典型例子，而柱面、锥面是 $K=0$ 的可展直纹面。}{For a ruled surface $\br(u,v)=\ba(u)+v\,\bb(u)$ the Gaussian curvature can be written as the displayed expression, which is nonpositive. Hence a ruled surface other than a developable one is a negatively curved surface; $K\equiv0$ exactly when $\det(\ba',\bb,\bb')=0$, that is when the surface is developable. The hyperbolic paraboloid (saddle) and the hyperboloid of one sheet are standard negative-curvature examples, while cylinders and cones are developable ruled surfaces with $K=0$.}
\APS{极小曲面是满足 $H\equiv0$ 的曲面，即两个主曲率处处互为相反数。悬链面（旋转面）与螺旋面是最经典的两个极小曲面，它们可以连续形变成彼此而保持极小性。肥皂膜在表面张力作用下使面积局部极小，其 Euler--Lagrange 方程正是 $H=0$；因此极小曲面理论是变分法与几何的交汇点。}{A minimal surface is one with $H\equiv0$, that is, the two principal curvatures are opposite at every point. The catenoid (a surface of revolution) and the helicoid are the two classical minimal surfaces, and they can be continuously deformed into each other while remaining minimal. A soap film minimizes area locally under surface tension, and its Euler--Lagrange equation is exactly $H=0$; minimal surface theory is therefore a meeting point of the calculus of variations and geometry.}

\subsection{球面、环面与悬链面的完整计算 | Complete computations for the sphere, torus and catenoid}
\EX{球面 $\br(u,v)=R(\sin u\cos v,\sin u\sin v,\cos u)$，外法向 $\bn=\br/R$。已求得 $E=R^2$、$F=0$、$G=R^2\sin^2u$。再求二阶导数：
\[
\br_{uu}=-R\bn,\qquad
\br_{uv}=R\cos u(-\sin v,\cos v,0),
\]
\[
\br_{vv}=-R\sin u(\sin u\cos v,\sin u\sin v,0)=-R\sin^2u\,\bn .
\]
故 $L=-R$、$M=0$、$N=-R\sin^2u$，于是 $K=(LN-M^2)/(EG-F^2)=R^2\sin^2u/(R^4\sin^2u)=1/R^2$，
\[
H=\frac{EN-2FM+GL}{2(EG-F^2)}=\frac{-R^3\sin^2u-R^3\sin^2u}{2R^4\sin^2u}=-\frac1R .
\]
两个主曲率都是 $-1/R$，与"球面处处是脐点"一致；$K>0$ 说明球面是椭圆点曲面，且弯曲程度随 $R$ 增大而减小。}{Sphere $\br(u,v)=R(\sin u\cos v,\sin u\sin v,\cos u)$ with outward normal $\bn=\br/R$. We already have $E=R^2$, $F=0$, $G=R^2\sin^2u$. The displayed second derivatives give $L=-R$, $M=0$, $N=-R\sin^2u$, hence $K=1/R^2$ and the displayed mean curvature $H=-1/R$. Both principal curvatures equal $-1/R$, consistent with the fact that every point of a sphere is umbilical; $K>0$ says the sphere consists of elliptic points, and the bending decreases as $R$ grows.}
\EX{环面 $\br(u,v)=((R+r\cos u)\cos v,(R+r\cos u)\sin v,r\sin u)$，$R>r>0$，取法向 $\bn=-(\cos u\cos v,\cos u\sin v,\sin u)$（此定向使 $H>0$；换成 $-\bn$ 则 $H$ 变号）。已求得 $E=r^2$、$F=0$、$G=(R+r\cos u)^2$。计算得 $\br_{uu}=-r(\cos u\cos v,\cos u\sin v,\sin u)=r\bn$，故 $L=r$；$\br_{uv}=r\sin u(\sin v,-\cos v,0)$ 与 $\bn$ 正交，故 $M=0$；
\[
\br_{vv}=-(R+r\cos u)(\cos v,\sin v,0),\qquad
N=(R+r\cos u)\cos u .
\]
于是
\[
K=\frac{LN}{EG}=\frac{\cos u}{r(R+r\cos u)},\qquad
H=\frac{R+2r\cos u}{2r(R+r\cos u)} .
\]
当 $u=0$（外赤道）$K=1/(r(R+r))>0$；当 $u=\pi$（内赤道）$K=-1/(r(R-r))<0$；当 $\cos u=0$（上下两圈）$K=0$。所以环面上同时出现椭圆点、双曲点和抛物点，是展示 $K$ 变号的典型曲面。}{Torus $\br(u,v)=((R+r\cos u)\cos v,(R+r\cos u)\sin v,r\sin u)$ with $R>r>0$ and the normal $\bn=-(\cos u\cos v,\cos u\sin v,\sin u)$, an orientation that makes $H>0$ (replacing $\bn$ by $-\bn$ changes the sign of $H$). We already have $E=r^2$, $F=0$, $G=(R+r\cos u)^2$. One computes $\br_{uu}=r\bn$, so $L=r$; $\br_{uv}$ is orthogonal to $\bn$, so $M=0$; and the displayed $\br_{vv}$ gives $N=(R+r\cos u)\cos u$. Hence the displayed $K$ and $H$. On the outer equator $u=0$, $K=1/(r(R+r))>0$; on the inner equator $u=\pi$, $K=-1/(r(R-r))<0$; on the two circles $\cos u=0$, $K=0$. The torus thus exhibits elliptic, hyperbolic and parabolic points at once and is the standard example of a surface where $K$ changes sign.}
\EX{悬链面 $\br(u,v)=(a\cosh(u/a)\cos v,a\cosh(u/a)\sin v,u)$，$a>0$。计算得
\[
E=\cosh^2\frac ua,\qquad F=0,\qquad G=a^2\cosh^2\frac ua,\qquad
I=\cosh^2\frac ua\left(\dd u^2+a^2\dd v^2\right).
\]
单位法向量为 $\bn=(-\cos v,-\sin v,\sinh(u/a))/\cosh(u/a)$。二阶导数为
\[
\br_{uu}=\frac1a\cosh\frac ua(\cos v,\sin v,0),\qquad
\br_{uv}=\sinh\frac ua(-\sin v,\cos v,0),\qquad
\br_{vv}=-a\cosh\frac ua(\cos v,\sin v,0),
\]
故 $L=-1/a$、$M=0$、$N=a$，从而
\[
K=\frac{LN}{EG}=-\frac{1}{a^2\cosh^4(u/a)},\qquad
H=\frac{EN+GL}{2(EG)}=\frac{a\cosh^2(u/a)-a\cosh^2(u/a)}{2a^2\cosh^4(u/a)}=0 .
\]
悬链面是极小曲面，$K<0$ 处处成立，且 $|K|$ 在腰部 $u=0$ 最大、向两端指数衰减。}{Catenoid $\br(u,v)=(a\cosh(u/a)\cos v,a\cosh(u/a)\sin v,u)$ with $a>0$. One computes the displayed first fundamental form, unit normal, second derivatives, and hence $L=-1/a$, $M=0$, $N=a$. Therefore $K=-1/(a^2\cosh^4(u/a))$ and $H=0$. The catenoid is a minimal surface, $K<0$ everywhere, and $|K|$ is largest at the waist $u=0$ and decays exponentially toward the ends.}
\ASK{第二基本形式与第一基本形式的区别是什么？}{What is the difference between the first and second fundamental forms?}{第一基本形式是切向量的内积，只度量曲面上的长度、角度和面积，是内蕴的；第二基本形式是曲线加速度的法向分量，描述曲面如何弯曲，依赖曲面在空间中的嵌入方式，是外蕴的（差一个符号）。二者合起来 $(I,II)$ 才在刚体运动下决定曲面。}{The first fundamental form is the inner product of tangent vectors and measures only length, angle and area on the surface, so it is intrinsic; the second fundamental form is the normal component of the acceleration of curves and describes how the surface bends, depending on the embedding in space, so it is extrinsic up to sign. Only the pair $(I,II)$ determines the surface up to rigid motion.}
\ASK{为什么球面上每一点都是脐点？}{Why is every point of a sphere umbilical?}{因为球面的单位法向量是位置向量除以半径，沿任意切方向法向量的变化率都相同：$\dd\bn(\bv)=-\bv/R$，所以 Weingarten 映射是数乘 $-1/R$，两个主曲率都是 $-1/R$，所有方向等价。这也是球面具有最大对称性的体现。}{Because the unit normal of a sphere is its position vector divided by the radius, the normal changes at the same rate in every tangent direction: $\dd\bn(\bv)=-\bv/R$. Hence the Weingarten map is the scalar $-1/R$, both principal curvatures are $-1/R$, and all directions are equivalent, which reflects the maximal symmetry of the sphere.}

\section[曲面论基本定理]{曲面论基本定理与绝妙定理 | The fundamental theorem and the Theorema Egregium}
\subsection{自然标架的运动公式 | Motion formulas of the natural frame}
\APS{曲面上每点有一个自然标架 $\{\br_u,\br_v,\bn\}$。把它们对 $u,v$ 求导并仍用这组标架展开，就得到 Gauss 公式与 Weingarten 公式：
\[
\br_{uu}=\Gamma^1_{11}\br_u+\Gamma^2_{11}\br_v+L\bn,\quad
\br_{uv}=\Gamma^1_{12}\br_u+\Gamma^2_{12}\br_v+M\bn,\quad
\br_{vv}=\Gamma^1_{22}\br_u+\Gamma^2_{22}\br_v+N\bn,
\]
\[
\bn_u=-h^1_1\br_u-h^2_1\br_v,\qquad
\bn_v=-h^1_2\br_u-h^2_2\br_v,
\]
其中 $L,M,N$ 是第二基本形式系数，$h^i_j$ 是 Weingarten 映射的矩阵元素。}{Every point of a surface carries the natural frame $\{\br_u,\br_v,\bn\}$. Differentiating these vectors with respect to $u,v$ and expanding again in the natural frame gives the Gauss and Weingarten formulas displayed, where $L,M,N$ are the coefficients of the second fundamental form and $h^i_j$ are the matrix entries of the Weingarten map.}
\APS{Christoffel 符号 $\Gamma^k_{ij}$ 完全由第一基本形式决定。由度量相容性 $\br_{ij}\cdot\br_k=\frac12(\partial_i E_{jk}+\partial_j E_{ik}-\partial_k E_{ij}$ 等关系可解出六个系数；例如在正交参数网 $F=0$ 下，
\[
\Gamma^1_{11}=\frac{E_u}{2E},\quad \Gamma^1_{12}=\frac{E_v}{2E},\quad \Gamma^1_{22}=-\frac{G_u}{2E},\quad
\Gamma^2_{11}=-\frac{E_v}{2G},\quad \Gamma^2_{12}=\frac{G_u}{2G},\quad \Gamma^2_{22}=\frac{G_v}{2G}.
\]
这些公式说明：自然标架的运动完全由 $E,F,G,L,M,N$ 这六个函数控制。}{The Christoffel symbols $\Gamma^k_{ij}$ are determined entirely by the first fundamental form. Metric compatibility, together with relations such as $\br_{ij}\cdot\br_k=\frac12(\partial_i E_{jk}+\partial_j E_{ik}-\partial_k E_{ij})$, determines the six coefficients; for example, for an orthogonal net $F=0$ they take the displayed values. These formulas show that the motion of the natural frame is controlled completely by the six functions $E,F,G,L,M,N$.}
\APS{运动公式的价值在于把"曲面的几何"翻译成"系数矩阵的微分方程"。若把 $\{\br_u,\br_v,\bn\}$ 的导数写成九个系数的线性组合，则这些系数由 $I,II$ 给出；反过来，知道这些系数就能通过解 ODE 重建标架，再积分得到曲面。这条路线的可积性条件正是下一节的 Gauss--Codazzi 方程。}{The value of the motion formulas is that they translate the geometry of a surface into differential equations for a coefficient matrix. Writing the derivatives of $\{\br_u,\br_v,\bn\}$ as linear combinations with nine coefficients, these coefficients are given by $I$ and $II$; conversely, knowing them one reconstructs the frame by solving ODEs and then integrates to obtain the surface. The integrability conditions of this programme are exactly the Gauss--Codazzi equations of the next subsection.}

\subsection{曲面的唯一性定理 | The uniqueness theorem for surfaces}
\APS{曲面论唯一性定理（Bonnet）：设两张正则曲面定义在同一参数区域上，且在对应点处有相同的第一、第二基本形式系数 $E,F,G,L,M,N$，则它们只相差一个刚体运动。换句话说，六个函数 $(E,F,G,L,M,N)$ 在刚体运动下完全决定曲面。}{The uniqueness theorem for surfaces (Bonnet): if two regular surfaces are defined on the same parameter domain and have the same coefficients $E,F,G,L,M,N$ of the first and second fundamental forms at corresponding points, then they differ only by a rigid motion. In other words, the six functions $(E,F,G,L,M,N)$ determine the surface completely up to rigid motion.}
\APS{证明思路与曲线论基本定理相同：在对应点把两个自然标架对齐，运动公式说明标架矩阵满足同一个一阶线性 ODE 组；两个解在初始时刻相同，由 Picard--Lindelöf 唯一性得标架处处相同，再积分比较位置向量，即得刚体运动。这里的关键是运动公式的系数只依赖 $I,II$，因此两个曲面共享同一个 ODE。}{The proof follows the same idea as the fundamental theorem of curve theory: align the two natural frames at corresponding points; the motion formulas say that the frame matrices satisfy the same first-order linear ODE system. Two solutions agree initially, so Picard--Lindelöf uniqueness makes the frames agree everywhere; integrating and comparing position vectors then gives a rigid motion. The key point is that the coefficients in the motion formulas depend only on $I,II$, so the two surfaces share the same ODE.}

\subsection{曲面论基本方程 | The fundamental equations: Gauss and Codazzi}
\APS{运动公式的系数不能任意给定，因为混合偏导数必须交换次序：$(\br_u)_{uv}=(\br_u)_{vu}$、$(\bn)_{uv}=(\bn)_{vu}$ 等。把这些等式用运动公式展开，比较 $\br_u,\br_v,\bn$ 的系数，就得到曲面论基本方程：一个 Gauss 方程和两个 Codazzi 方程。}{The coefficients in the motion formulas cannot be prescribed arbitrarily, because mixed partial derivatives must commute: $(\br_u)_{uv}=(\br_u)_{vu}$, $(\bn)_{uv}=(\bn)_{vu}$, and so on. Expanding these identities with the motion formulas and comparing coefficients of $\br_u,\br_v,\bn$ gives the fundamental equations of surface theory: one Gauss equation and two Codazzi equations.}
\FORMULA{Gauss 方程与 Codazzi 方程}{the Gauss and Codazzi equations}
\[
K=\frac{LN-M^2}{EG-F^2},
\]
\[
L_v-M_u=\Gamma^1_{12}L+\left(\Gamma^2_{12}-\Gamma^1_{11}\right)M-\Gamma^2_{11}N,
\]
\[
N_u-M_v=-\Gamma^1_{22}L+\left(\Gamma^1_{12}-\Gamma^2_{22}\right)M+\Gamma^2_{12}N .
\]
\APS{注意下标约定：本文用 $L_v=\partial L/\partial v$、$M_u=\partial M/\partial u$。Gauss 方程把 Gauss 曲率与第二基本形式联系起来；Codazzi 方程描述第二基本形式如何随参数变化。三个方程中只有三个是独立的（Gauss 方程可以看作一个标量方程，Codazzi 两个），它们正是六个函数 $E,F,G,L,M,N$ 必须满足的全部相容性条件。}{Note the subscript convention: here $L_v=\partial L/\partial v$ and $M_u=\partial M/\partial u$. The Gauss equation links the Gaussian curvature to the second fundamental form, and the Codazzi equations describe how the second fundamental form varies with the parameters. Only three of the equations are independent (the Gauss equation is one scalar equation and the Codazzi equations are two), and they are exactly the complete compatibility conditions that the six functions $E,F,G,L,M,N$ must satisfy.}
\APS{基本方程的地位可以这样理解：它们是唯一性与存在性定理的分界线。满足 Gauss--Codazzi 方程的六个函数一定来自某张曲面（存在性），而任何曲面的六个系数一定满足这些方程（必要性）；因此这组方程恰好刻画了什么样的 $I,II$ 能成为曲面。}{The role of the fundamental equations can be understood as follows: they are the dividing line between the uniqueness and existence theorems. Six functions satisfying the Gauss--Codazzi equations always come from some surface (existence), and the six coefficients of any surface always satisfy these equations (necessity); the equations therefore characterize exactly which pairs $I,II$ can arise from a surface.}

\subsection{曲面的存在性定理 | The existence theorem for surfaces}
\APS{存在性定理是唯一性定理的逆命题：给定定义在单连通参数区域上的六个 $C^2$ 函数 $E,F,G$ 与 $L,M,N$，若 $E>0$、$G>0$、$EG-F^2>0$（正定），且它们满足 Gauss--Codazzi 方程，则存在正则曲面以它们为第一、第二基本形式，并且这样的曲面在刚体运动下唯一。}{The existence theorem is the converse of the uniqueness theorem: given six $C^2$ functions $E,F,G$ and $L,M,N$ on a simply connected parameter domain, if $E>0$, $G>0$, $EG-F^2>0$ (positive definite) and they satisfy the Gauss--Codazzi equations, then there is a regular surface having them as first and second fundamental forms, and such a surface is unique up to rigid motion.}
\APS{证明思路：把运动公式看作关于标架矩阵 $\{\br_u,\br_v,\bn\}$ 的一阶线性 ODE 组，任取一组初始标架（正交且右手）解之；Gauss--Codazzi 方程保证解出的标架自动保持正交，且由 $\br_{uv}=\br_{vu}$ 保证存在向量函数 $\br(u,v)$ 使 $\br_u,\br_v$ 为解中的前两个向量；最后 $\br$ 的法向量就是解中的第三个向量。这一步用到 Frobenius 可积性定理，是"用 ODE 构造几何对象"的典范。}{The proof idea: regard the motion formulas as a first-order linear ODE system for the frame matrix $\{\br_u,\br_v,\bn\}$ and solve it with any orthonormal right-handed initial frame; the Gauss--Codazzi equations guarantee that the solution frame stays orthonormal, and the symmetry $\br_{uv}=\br_{vu}$ guarantees the existence of a vector function $\br(u,v)$ whose partial derivatives are the first two vectors of the solution; finally the normal of $\br$ is the third vector. This step uses the Frobenius integrability theorem and is a model of constructing a geometric object by ODEs.}

\subsection{Gauss 绝妙定理 | Gauss's Theorema Egregium}
\APS{Gauss 绝妙定理：Gauss 曲率 $K$ 只由第一基本形式 $E,F,G$ 决定，与第二基本形式无关。换句话说，$K$ 是内蕴量；两个曲面若保长对应，则对应点的 $K$ 相等。Gauss 称其为"绝妙定理"，因为它说明一个看似外蕴的量（法向量变化的速率）其实完全由曲面上的度量决定。}{Gauss's Theorema Egregium: the Gaussian curvature $K$ is determined by the first fundamental form $E,F,G$ alone, independently of the second fundamental form. In other words, $K$ is intrinsic: if two surfaces are isometric then their Gaussian curvatures agree at corresponding points. Gauss called it the "remarkable theorem" because it says that a seemingly extrinsic quantity, the rate of change of the normal, is in fact determined entirely by the metric on the surface.}
\APS{证明的关键是把 Gauss 方程 $K=(LN-M^2)/(EG-F^2)$ 的右端用第一基本形式表示。最直接的形式是对正交参数网：
\[
K=-\frac{1}{2\sqrt{EG}}\left[\frac{\partial}{\partial v}\!\left(\frac{E_v}{\sqrt{EG}}\right)
+\frac{\partial}{\partial u}\!\left(\frac{G_u}{\sqrt{EG}}\right)\right].
\]
这个公式只含 $E,G$ 及其一、二阶导数，因此 $K$ 是内蕴的。验证：球面 $E=R^2$、$G=R^2\sin^2u$ 代入得 $K=1/R^2$；柱面 $E=R^2$、$G=1$ 代入得 $K=0$，与直接计算一致。}{The key to the proof is to express the right-hand side of the Gauss equation $K=(LN-M^2)/(EG-F^2)$ through the first fundamental form. The most direct form is for an orthogonal net, displayed; it contains only $E,G$ and their first and second derivatives, so $K$ is intrinsic. Check: substituting the sphere $E=R^2$, $G=R^2\sin^2u$ gives $K=1/R^2$, and the cylinder $E=R^2$, $G=1$ gives $K=0$, in agreement with direct computation.}
\APS{绝妙定理最重要的推论是"曲率是等距不变量"。平面与柱面都是 $K=0$ 的可展曲面，可以互相等距展开；球面 $K=1/R^2>0$，不可能与平面等距，因此任何地图投影都必须在长度或角度上妥协：Mercator 保角但面积失真，等积投影保面积但角度失真。这就是"球面不能摊平"的精确数学表述。}{The most important corollary of the Theorema Egregium is that curvature is an isometric invariant. The plane and the cylinder are both developable surfaces with $K=0$ and can be unrolled into each other isometrically; the sphere has $K=1/R^2>0$ and cannot be isometric to a plane, so every map projection must distort either length or angle: Mercator is conformal but distorts area, and equal-area projections preserve area but distort angles. This is the precise mathematical statement that a sphere cannot be flattened.}
\APS{绝妙定理也解释了 Gauss 映射的几何意义：$\bn:S\to S^2$ 把曲面映到单位球面，$K$ 恰好是这一映射的 Jacobi 行列式，$\int_S|K|\,\dd A$ 是球面像的面积。$K$ 的符号对应映射保向或反向，这与椭圆点、双曲点的分类一致。第 6 章的 Gauss--Bonnet 公式将把 $K$ 的积分与曲面的整体拓扑联系起来，可以看作绝妙定理的全局版本。}{The Theorema Egregium also explains the geometry of the Gauss map: $\bn:S\to S^2$ maps the surface to the unit sphere, $K$ is exactly the Jacobian determinant of this map, and $\int_S|K|\,\dd A$ is the area of the spherical image. The sign of $K$ records whether the map preserves or reverses orientation, consistent with the classification into elliptic and hyperbolic points. The Gauss--Bonnet formula of Chapter 6 links the integral of $K$ to the global topology of the surface and may be regarded as the global version of the Theorema Egregium.}
\ASK{为什么说绝妙定理是"绝妙"的？}{Why is the Theorema Egregium called remarkable?}{因为 $K$ 最初由法向量的变化（第二基本形式）定义，看起来依赖曲面在空间中的嵌入；定理却说 $K$ 只由第一基本形式决定，因而在弯曲曲面（只要不拉伸）时保持不变。这解释了为什么地图不能既保长又保角地表示球面，也把内蕴几何从外蕴描述中分离出来。}{Because $K$ was originally defined through the change of the normal (the second fundamental form) and seems to depend on how the surface is embedded in space; the theorem says instead that $K$ depends only on the first fundamental form and is therefore preserved under bending without stretching. This explains why no map can represent a sphere both length-preserving and angle-preserving, and it separates intrinsic geometry from extrinsic description.}
\ASK{Gauss--Codazzi 方程的作用是什么？}{What is the role of the Gauss--Codazzi equations?}{它们是六个系数 $E,F,G,L,M,N$ 的相容性条件，等价于曲面存在且唯一（到刚体运动）的充要条件。Gauss 方程把 $K$ 与 $II$ 联系起来，Codazzi 方程保证第二基本形式的变化与第一基本形式匹配；没有它们，就无法由六个函数构造出曲面。}{They are the compatibility conditions on the six coefficients $E,F,G,L,M,N$ and are equivalent to the necessary and sufficient conditions for a surface to exist and be unique up to rigid motion. The Gauss equation links $K$ to $II$, and the Codazzi equations ensure that the variation of the second fundamental form matches the first; without them the six functions cannot be assembled into a surface.}

\section[测地几何]{测地曲率、测地线与 Gauss--Bonnet | Geodesic curvature, geodesics and Gauss--Bonnet}
\subsection{测地曲率与测地挠率 | Geodesic curvature and geodesic torsion}
\APS{曲面上一条曲线有两种弯曲：沿曲面"转弯"和离开曲面"抬头"。设曲线单位速度、Frenet 标架为 $\{\bT,\bN,\bB\}$，曲面单位法向量为 $\bn$。把加速度 $\bT'$ 分解为法向与切向两部分：
\[
\bT'=\kappa_n\bn+\kappa_g(\bn\times\bT),\qquad
\kappa_n=II(\bT,\bT),\qquad
\kappa_g=\bT'\cdot(\bn\times\bT),
\]
其中 $\kappa_n$ 是法曲率，$\kappa_g$ 是测地曲率。由 $\bT'\perp\bT$ 得 $\kappa^2=\kappa_n^2+\kappa_g^2$：曲线的总曲率由"曲面自身的弯曲"与"曲线在曲面内转弯"两部分组成。}{A curve on a surface bends in two ways: it turns within the surface and it tilts away from the surface. Let the curve have unit speed and Frenet frame $\{\bT,\bN,\bB\}$, and let $\bn$ be the unit normal of the surface. Decompose the acceleration $\bT'$ into normal and tangential parts as displayed, where $\kappa_n$ is the normal curvature and $\kappa_g$ the geodesic curvature. Since $\bT'\perp\bT$, one has $\kappa^2=\kappa_n^2+\kappa_g^2$: the total curvature of the curve splits into the bending of the surface itself and the turning of the curve inside the surface.}
\APS{在正交参数网 $F=0$ 下，测地曲率有便于计算的 Liouville 公式。设曲线切向量与 $u$ 曲线的夹角为 $\theta$，则
\[
\kappa_g=\frac{\dd\theta}{\dd s}
+\frac{1}{2\sqrt E}\frac{\partial\ln G}{\partial u}\sin\theta
-\frac{1}{2\sqrt G}\frac{\partial\ln E}{\partial v}\cos\theta .
\]
这个公式把"曲线在曲面内转弯"分成两项：$\dd\theta/\dd s$ 是切向量相对于参数网转动的速率，其余两项来自参数网自身的非正交性与尺度变化。}{For an orthogonal net $F=0$ there is a convenient Liouville formula for the geodesic curvature. If $\theta$ is the angle between the tangent of the curve and the $u$-curve, then the displayed formula holds. It splits the turning of the curve inside the surface into two parts: $\dd\theta/\dd s$ is the rate at which the tangent turns relative to the parameter net, and the remaining terms come from the non-orthogonality and the changing scale of the net.}
\APS{类似地定义测地挠率 $\tau_g=-\bn'\cdot(\bn\times\bT)$。它与曲线的挠率 $\tau$ 以及曲面法向量与曲线主法向量的夹角 $\varphi$ 满足 $\tau_g=\tau+\dd\varphi/\dd s$；若用主方向计算，则
\[
\tau_g=\frac12(\kappa_1-\kappa_2)\sin2\theta ,
\]
其中 $\theta$ 是切向量与第一主方向的夹角。因此沿主方向（曲率线）测地挠率为零；这也是"曲率线"在曲面上的另一刻画。}{Similarly define the geodesic torsion $\tau_g=-\bn'\cdot(\bn\times\bT)$. It satisfies $\tau_g=\tau+\dd\varphi/\dd s$, where $\tau$ is the torsion of the curve and $\varphi$ the angle between the surface normal and the principal normal of the curve; in terms of the principal directions, the displayed formula holds, where $\theta$ is the angle between the tangent and the first principal direction. Hence the geodesic torsion vanishes along principal directions (lines of curvature), which gives another characterization of lines of curvature on a surface.}
\EX{球面 $R=1$ 上纬线 $u=u_0$ 的测地曲率。纬线是 $v$ 曲线，$E=1$、$G=\sin^2u$，$\theta=\pi/2$，故 Liouville 公式给 $\kappa_g=\dfrac{1}{2\sqrt E}\dfrac{\partial\ln G}{\partial u}\sin\theta=\dfrac12\cdot 2\cot u_0\cdot 1=\cot u_0$。赤道 $u_0=\pi/2$ 处 $\kappa_g=0$，纬线是测地线；其余纬线 $\kappa_g\ne0$，不是测地线，且越靠近极点弯曲越厉害。}{Geodesic curvature of the latitude circle $u=u_0$ on the unit sphere. The latitude is a $v$-curve with $E=1$, $G=\sin^2u$ and $\theta=\pi/2$, so the Liouville formula gives $\kappa_g=(1/(2\sqrt E))(\partial_u\ln G)\sin\theta=(1/2)\cdot2\cot u_0=\cot u_0$. At the equator $u_0=\pi/2$ one has $\kappa_g=0$ and the latitude is a geodesic; other latitudes have $\kappa_g\ne0$ and are not geodesics, bending more strongly near the poles.}

\subsection{测地线 | Geodesics}
\APS{测地线是测地曲率恒为零的曲线：$\kappa_g\equiv0$，等价地 $\bT'$ 处处平行于曲面法向量 $\bn$，即曲线在曲面内不"转弯"，其加速度只改变它离开曲面的程度。测地线是曲面上"最直的线"，也是局部最短线：任一点附近，连接两点的最短曲线必是测地线；反之测地线在小范围内实现最短距离，但整体上不一定（球面上大圆的两段弧一长一短）。}{A geodesic is a curve with identically zero geodesic curvature, $\kappa_g\equiv0$; equivalently $\bT'$ is everywhere parallel to the surface normal $\bn$, so the curve does not turn inside the surface and its acceleration only changes how it leaves the surface. Geodesics are the straightest curves on a surface and locally the shortest: near any point the shortest curve joining two points is a geodesic, and conversely a geodesic realizes the shortest distance in a small range, though not globally (the two arcs of a great circle on a sphere have different lengths).}
\APS{测地线是长度泛函的临界点，也可以由能量泛函的 Euler--Lagrange 方程得到。设曲线用参数 $t$ 给出，能量为 $\frac12\int(Eu'^2+2Fu'v'+Gv'^2)\dd t$，其 Euler--Lagrange 方程就是测地线方程
\[
u''+\Gamma^1_{11}(u')^2+2\Gamma^1_{12}u'v'+\Gamma^1_{22}(v')^2=0,\qquad
v''+\Gamma^2_{11}(u')^2+2\Gamma^2_{12}u'v'+\Gamma^2_{22}(v')^2=0 .
\]
这是二阶拟线性 ODE 组，给定一点及一个单位切方向，由解的存在唯一性得到唯一测地线；所有测地线构成指数映射 $\exp_p:T_pS\to S$，这是测地坐标与法坐标的基础。}{Geodesics are critical points of the length functional and can also be obtained from the Euler--Lagrange equations of the energy functional. For a curve with parameter $t$, the energy is $\frac12\int(Eu'^2+2Fu'v'+Gv'^2)\dd t$, and its Euler--Lagrange equations are the displayed geodesic equations. This is a second-order quasilinear ODE system; given a point and a unit tangent direction, existence and uniqueness of solutions give a unique geodesic, and all geodesics define the exponential map $\exp_p:T_pS\to S$, the basis of geodesic and normal coordinates.}
\APS{球面的测地线方程最容易写出来。取 $E=R^2$、$G=R^2\sin^2u$，忽略 $R$ 的常数因子得
\[
u''-\sin u\cos u\,(v')^2=0,\qquad v''+2\cot u\,u'v'=0 .
\]
它们的解是大圆：过球心的平面与球面的交线。从对称性看，大圆上每点的加速度都指向球心，即平行于法向量，故测地曲率为零；反过来，球面上任一测地线必为大圆的一段。}{The geodesic equations of a sphere are the easiest to write down. With $E=R^2$ and $G=R^2\sin^2u$, dropping the constant factor of $R$ gives the displayed equations. Their solutions are great circles, the intersections of the sphere with planes through its center. By symmetry the acceleration at each point of a great circle points to the center, hence is parallel to the normal, so the geodesic curvature is zero; conversely every geodesic on a sphere is an arc of a great circle.}
\APS{柱面上的测地线可以完全不靠计算理解：柱面可以等距展开成平面，等距映射把测地线映为测地线（因为测地曲率是内蕴量）。平面上的测地线是直线，展开回去就得到柱面上的螺旋线（包括母线、平行圆和一般螺线）。这个例子说明"内蕴几何 + 展开"是研究可展曲面上测地线的最快方法。}{Geodesics on a cylinder can be understood without any computation: the cylinder can be unrolled isometrically onto a plane, and an isometry maps geodesics to geodesics because geodesic curvature is intrinsic. Geodesics in the plane are straight lines, so unrolling back gives helices on the cylinder (including generators, parallel circles and general helices). This example shows that "intrinsic geometry plus development" is the fastest way to study geodesics on developable surfaces.}

\subsection{测地坐标系与法坐标系 | Geodesic and normal coordinates}
\APS{以点 $p$ 为中心，用测地极坐标 $(r,\theta)$ 表示曲面上的点：$r$ 是从 $p$ 出发的测地线长度，$\theta$ 是初始切方向的角度。此时第一基本形式化为
\[
I=\dd r^2+G(r,\theta)\,\dd\theta^2,
\]
其中 $G(0,\theta)=0$、$G_r(0,\theta)=1$，且 $G_{rr}(0,\theta)=0$。这就是 Gauss 引理：从 $p$ 出发的测地线与以 $p$ 为心的测地圆正交，因此 $F=0$、$E=1$。}{With a point $p$ as center, use geodesic polar coordinates $(r,\theta)$ to parametrize the surface: $r$ is the length of a geodesic starting at $p$ and $\theta$ is the angle of the initial tangent direction. The first fundamental form then becomes the displayed expression, where $G(0,\theta)=0$, $G_r(0,\theta)=1$ and $G_{rr}(0,\theta)=0$. This is Gauss's lemma: geodesics from $p$ are orthogonal to the geodesic circles centered at $p$, so $F=0$ and $E=1$.}
\APS{由 $\sqrt G$ 在 $r=0$ 附近的展开可以读出曲率。可以证明
\[
G(r,\theta)=r^2-\frac{K(p)}{3}r^4+O(r^5),
\]
于是以 $p$ 为心、半径 $r$ 的测地圆周长为 $2\pi r\left(1-\frac{K(p)}{6}r^2+O(r^4)\right)$，面积为 $\pi r^2\left(1-\frac{K(p)}{12}r^2+O(r^4)\right)$。当 $K>0$ 时圆周比 Euclidean 圆短、面积比 Euclidean 小（球面上经线向极点汇聚）；当 $K<0$ 时相反。测地坐标因此是"用测量长度来感知曲率"的工具。}{The expansion of $\sqrt G$ near $r=0$ reveals the curvature. One can show the displayed expansion of $G(r,\theta)$, hence the circumference of the geodesic circle of radius $r$ centered at $p$ is $2\pi r(1-\frac{K(p)}{6}r^2+O(r^4))$ and its area is $\pi r^2(1-\frac{K(p)}{12}r^2+O(r^4))$. When $K>0$ the circumference is shorter and the area smaller than in the Euclidean plane (meridians converge toward the poles on a sphere), and the opposite holds when $K<0$. Geodesic coordinates are therefore the tool for perceiving curvature by measuring lengths.}
\APS{法坐标系是测地极坐标的直角版本：取 $T_pS$ 的一组单位正交基 $\be_1,\be_2$，令 $\bx=x^1\be_1+x^2\be_2$，定义 $\exp_p(\bx)$ 为从 $p$ 出发、初始切向量为 $\bx$ 的测地线在参数 $1$ 处的点。在法坐标下 $g_{ij}(p)=\delta_{ij}$、$\partial_k g_{ij}(p)=0$，即度量在 $p$ 处与 Euclidean 度量有二阶接近；这就是"任意曲面局部近似于平面"的精确说法。}{Normal coordinates are the Cartesian version of geodesic polar coordinates: choose an orthonormal basis $\be_1,\be_2$ of $T_pS$, write $\bx=x^1\be_1+x^2\be_2$, and define $\exp_p(\bx)$ as the point at parameter $1$ on the geodesic starting at $p$ with initial tangent vector $\bx$. In normal coordinates $g_{ij}(p)=\delta_{ij}$ and $\partial_k g_{ij}(p)=0$, so the metric agrees with the Euclidean metric at $p$ to second order; this is the precise statement that every surface is locally approximated by a plane.}

\subsection{常曲率曲面 | Surfaces of constant curvature}
\APS{若 Gauss 曲率 $K$ 在曲面上为常数，称其为常曲率曲面。Minding 定理说：两张常曲率曲面若 $K$ 相同，则局部等距。因此常曲率曲面只有三种局部模型：$K=1/R^2>0$ 的球面（半径 $R$）、$K=0$ 的平面、$K=-1/a^2<0$ 的双曲平面。这正是"空间形式"在二维的情形。}{If the Gaussian curvature $K$ is constant on a surface, it is called a surface of constant curvature. Minding's theorem says that two surfaces of constant curvature with the same $K$ are locally isometric. Hence there are only three local models: the sphere of radius $R$ with $K=1/R^2>0$, the plane with $K=0$, and the hyperbolic plane with $K=-1/a^2<0$. These are the two-dimensional cases of the space forms.}
\APS{常曲率曲面上测地三角形的内角和由 $K$ 决定。对由测地线围成的三角形，Gauss--Bonnet 给
\[
\alpha+\beta+\gamma-\pi=\iint_\triangle K\,\dd A=K\cdot\operatorname{Area}(\triangle).
\]
因此 $K>0$ 时内角和大于 $\pi$（球面三角形），$K=0$ 时等于 $\pi$（平面三角形），$K<0$ 时小于 $\pi$（双曲三角形）。这是"曲率 = 角度亏损"的最直观表达，也是非欧几何的起点。}{On a surface of constant curvature the angle sum of a geodesic triangle is determined by $K$. For a triangle bounded by geodesics, Gauss--Bonnet gives the displayed identity, so the angle sum is greater than $\pi$ when $K>0$ (spherical triangles), equal to $\pi$ when $K=0$ (plane triangles), and less than $\pi$ when $K<0$ (hyperbolic triangles). This is the most intuitive expression of "curvature equals angle defect" and the starting point of non-Euclidean geometry.}

\subsection{平行移动与和乐 | Parallel transport and holonomy}
\APS{沿曲线 $\bgamma$ 的向量场 $V$ 若满足 $\nabla_{\bgamma'}V=0$，即 $V'$ 没有切向分量（$V'$ 平行于法向量 $\bn$），则称 $V$ 沿曲线平行移动。这个方程是一阶线性 ODE，给定初值 $V(0)$ 有唯一解，且解保持长度与夹角，因此平行移动给出切空间之间的等距。}{A vector field $V$ along a curve $\bgamma$ is said to be parallel transported if $\nabla_{\bgamma'}V=0$, that is if $V'$ has no tangential component ($V'$ is parallel to the normal $\bn$). This equation is a first-order linear ODE with a unique solution for given initial value $V(0)$, and the solution preserves lengths and angles, so parallel transport gives isometries between tangent spaces.}
\APS{平行移动一般依赖路径：沿闭曲线平行移动一周后，向量会转过一个角度，称为和乐（holonomy）。对单位速度曲线，若 $V$ 与切向量的夹角为 $\alpha$，则 $\dd\alpha/\dd s=-\kappa_g$；因此沿闭曲线一周的转角为 $-\oint\kappa_g\dd s$。由 Gauss--Bonnet，它等于 $-2\pi+\iint_DK\dd A$，模 $2\pi$ 后就是 $\iint_DK\dd A$：和乐角（模 $2\pi$）正是所围区域的总曲率。球面上沿纬线 $u=u_0$ 平移一周的转角为 $-2\pi\cos u_0$，其模 $2\pi$ 的值 $2\pi(1-\cos u_0)$ 等于所围球冠的面积，赤道上为零，越靠近极点越大。}{Parallel transport generally depends on the path: after transporting a vector once around a closed curve it returns rotated by an angle called the holonomy. For a unit-speed curve, if $\alpha$ is the angle between $V$ and the tangent, then $\dd\alpha/\dd s=-\kappa_g$; hence the rotation around a closed curve is $-\oint\kappa_g\dd s$. By Gauss--Bonnet this equals $-2\pi+\iint_DK\dd A$, which modulo $2\pi$ is $\iint_DK\dd A$: the holonomy angle modulo $2\pi$ is exactly the total curvature of the enclosed region. On a sphere, transport once around the latitude $u=u_0$ rotates by $-2\pi\cos u_0$, whose value modulo $2\pi$, namely $2\pi(1-\cos u_0)$, equals the area of the enclosed spherical cap; it is zero at the equator and grows toward the poles.}
\APS{和乐为零当且仅当所围区域的总曲率为零，这刻画了"平坦"的曲面片。平面与可展曲面（$K=0$）上平行移动与路径无关，可以整体定义平行向量场；球面与双曲面上则依赖路径。这个现象在物理中就是"弯曲空间中的方向不可积"，也是规范场论中和乐概念的几何原型。}{The holonomy vanishes exactly when the total curvature of the enclosed region is zero, which characterizes "flat" pieces of a surface. On the plane and on developable surfaces ($K=0$) parallel transport is path-independent and global parallel fields exist; on spheres and hyperbolic surfaces it depends on the path. In physics this phenomenon is the non-integrability of direction in curved space, and it is the geometric prototype of holonomy in gauge theory.}

\subsection{抽象曲面与内蕴几何 | Abstract surfaces and intrinsic geometry}
\APS{抽象曲面是不依赖 $\RR^3$ 的二维流形：它由坐标卡（图册）覆盖，每张卡把曲面的一部分映到平面区域，坐标变换是光滑的；再在每点给定一个正定二次型作为度量（第一基本形式），就得到抽象曲面。长度、角度、面积、Christoffel 符号、测地线、平行移动都可以只由度量定义，不需要外部空间。}{An abstract surface is a two-dimensional manifold that does not depend on $\RR^3$: it is covered by coordinate charts, each mapping a piece of the surface to a plane domain with smooth transition maps; a positive definite quadratic form (the first fundamental form) is then prescribed at every point as the metric. Length, angle, area, Christoffel symbols, geodesics and parallel transport can all be defined from the metric alone, without any ambient space.}
\APS{抽象曲面上的"几何学"就是内蕴几何。平面与柱面作为抽象曲面是等距的（有相同度量），尽管它们在 $\RR^3$ 中的形状完全不同；平坦环面、双曲平面、锥面（除顶点外）都是重要的例子。内蕴几何关心的是"住在曲面上的二维生物"能测量到的性质：它无法区分平面与柱面，但能通过测地三角形的内角和判断 $K$ 的符号。}{The "geometry" on an abstract surface is intrinsic geometry. The plane and the cylinder are isometric as abstract surfaces (they have the same metric) even though their shapes in $\RR^3$ are completely different; the flat torus, the hyperbolic plane and the cone with its vertex removed are important examples. Intrinsic geometry concerns what a two-dimensional being living on the surface could measure: it cannot distinguish a plane from a cylinder, but it can determine the sign of $K$ from the angle sums of geodesic triangles.}
\APS{抽象曲面上的曲率可以用 Gauss 方程内蕴地定义：由度量算出 Christoffel 符号，再定义 Riemann 曲率张量；在二维情形，Riemann 张量只有一个独立分量，它就等于 $K$。这正是绝妙定理的抽象版本：曲率不需要第二基本形式，也不需要嵌入，它是度量自身的性质。}{The curvature of an abstract surface can be defined intrinsically through the Gauss equation: compute the Christoffel symbols from the metric and then define the Riemann curvature tensor; in two dimensions the Riemann tensor has only one independent component, and it equals $K$. This is the abstract version of the Theorema Egregium: curvature needs neither the second fundamental form nor an embedding, but is a property of the metric itself.}

\subsection{Gauss--Bonnet 公式 | The Gauss--Bonnet formula}
\APS{局部 Gauss--Bonnet 公式把沿边界的测地曲率与内部的 Gauss 曲率联系起来。设 $C$ 是曲面上一条简单闭曲线，围成区域 $D$，则
\[
\oint_C\kappa_g\,\dd s+\iint_DK\,\dd A=2\pi .
\]
若 $C$ 有角点，设外角为 $\varepsilon_i$（即转向角），则公式改为
\[
\oint_C\kappa_g\,\dd s+\sum_i\varepsilon_i+\iint_DK\,\dd A=2\pi .
\]
特别地，测地三角形的三个外角之和为 $2\pi-\iint_\triangle K\dd A$，等价于 $\alpha+\beta+\gamma-\pi=\iint_\triangle K\dd A$。}{The local Gauss--Bonnet formula links the geodesic curvature along a boundary to the Gaussian curvature inside. Let $C$ be a simple closed curve on a surface bounding a region $D$; then the first displayed formula holds. If $C$ has corners with exterior angles $\varepsilon_i$ (turning angles), the formula becomes the second displayed one. In particular, for a geodesic triangle the sum of the three exterior angles is $2\pi-\iint_\triangle K\dd A$, equivalently $\alpha+\beta+\gamma-\pi=\iint_\triangle K\dd A$.}
\APS{整体 Gauss--Bonnet 公式把曲率积分与拓扑不变量联系起来。设 $M$ 是紧致、可定向的闭曲面，Euler 示性数为 $\chi(M)$，则
\[
\iint_MK\,\dd A=2\pi\chi(M).
\]
球面 $\chi=2$，故 $\iint_{S^2}K\dd A=4\pi$；环面 $\chi=0$，故总曲率为零；亏格 $g$ 的曲面 $\chi=2-2g$，总曲率为 $4\pi(1-g)$。这个公式说明：无论怎样弯曲曲面（只要不撕裂），总曲率是拓扑不变量，只由亏格决定。}{The global Gauss--Bonnet formula links the integral of curvature to a topological invariant. Let $M$ be a compact, orientable closed surface with Euler characteristic $\chi(M)$; then the displayed identity holds. The sphere has $\chi=2$, hence $\iint_{S^2}K\dd A=4\pi$; the torus has $\chi=0$, hence zero total curvature; a surface of genus $g$ has $\chi=2-2g$ and total curvature $4\pi(1-g)$. The formula says that however the surface is bent (without tearing), the total curvature is a topological invariant determined by the genus alone.}
\APS{Euler 示性数可以用剖分计算：$\chi=V-E+F$（顶点数减边数加面数）。Gauss--Bonnet 因此也适用于多面体：把每个面的角亏损（$2\pi$ 减去该顶点处各面角之和）相加，得到 $4\pi$，这是 Descartes 的角亏损定理，也是离散版本的 Gauss--Bonnet。对带边界的曲面，公式为 $\iint_MK\dd A+\oint_{\partial M}\kappa_g\dd s=2\pi\chi(M)$，其中边界按诱导定向取测地曲率。}{The Euler characteristic can be computed from a triangulation: $\chi=V-E+F$ (vertices minus edges plus faces). Gauss--Bonnet therefore also applies to polyhedra: summing the angle defect at each vertex ($2\pi$ minus the sum of the face angles there) gives $4\pi$, which is Descartes' angle-defect theorem and the discrete version of Gauss--Bonnet. For a surface with boundary the formula reads $\iint_MK\dd A+\oint_{\partial M}\kappa_g\dd s=2\pi\chi(M)$, with the geodesic curvature taken with respect to the induced orientation of the boundary.}
\EX{球面三角形与 Gauss--Bonnet 的验证。取单位球面 $R=1$ 上由三个大圆围成的卦限三角形：顶点为 $(1,0,0)$、$(0,1,0)$、$(0,0,1)$，三条边都是大圆弧。每两条边在顶点处的夹角都是 $\pi/2$，故内角和为 $3\pi/2$，超出 $\pi/2$；三角形面积是球面面积的 $1/8$，即 $4\pi/8=\pi/2$。而 $\iint_\triangle K\dd A=1\cdot\pi/2=\pi/2$，与角度超出量一致，Gauss--Bonnet 成立。}{Verification of Gauss--Bonnet with a spherical triangle. On the unit sphere take the octant triangle bounded by three great circles with vertices $(1,0,0)$, $(0,1,0)$ and $(0,0,1)$; each pair of sides meets at angle $\pi/2$, so the angle sum is $3\pi/2$, exceeding $\pi$ by $\pi/2$. The area of the triangle is $1/8$ of the sphere, that is $4\pi/8=\pi/2$. Since $\iint_\triangle K\dd A=1\cdot\pi/2=\pi/2$, the angle excess and the curvature integral agree, confirming Gauss--Bonnet.}
\EX{球面大圆与柱面螺线都是测地线。球面大圆：它在一个过球心的平面内，曲率向量指向球心，与球面法向量平行，故 $\kappa_g=0$；也可由球面的对称性直接看出。柱面螺线：把柱面等距展开成平面 $\br(u,v)\mapsto(u,v)$，螺线 $\bgamma(t)=(R\cos t,R\sin t,ht)$ 展开为直线 $(Rt,ht)$，直线是平面的测地线，而等距保持测地曲率为零，故螺线是柱面的测地线。这说明判断测地线时，"展开成平面"往往比解测地线方程更快。}{Great circles on a sphere and helices on a cylinder are geodesics. A great circle lies in a plane through the center, its curvature vector points to the center and is parallel to the surface normal, so $\kappa_g=0$; this is also clear from the symmetry of the sphere. For a helix on a cylinder, unroll the cylinder isometrically onto the plane by $\br(u,v)\mapsto(u,v)$; the helix $\bgamma(t)=(R\cos t,R\sin t,ht)$ unfolds to the straight line $(Rt,ht)$, which is a geodesic of the plane, and isometries preserve zero geodesic curvature, so the helix is a geodesic of the cylinder. This shows that for recognizing geodesics, "unrolling to a plane" is often faster than solving the geodesic equations.}
\ASK{测地线一定是曲面上两点间的最短线吗？}{Is a geodesic always the shortest path between two points?}{不一定。测地线是局部最短的：每点附近存在邻域使其实现最短距离；但整体上可能有另一条更短的曲线，或测地线自身绕远（球面大圆的两段弧一长一短）。闭测地线是周期解，如球面大圆和环面上的闭测地线，只满足"局部最短"或"变分为零"，不一定是全局最短。}{Not necessarily. A geodesic is locally shortest: near each point there is a neighbourhood in which it realizes the shortest distance; globally, however, another curve may be shorter, or the geodesic itself may take the long way around (the two arcs of a great circle have different lengths). Closed geodesics are periodic solutions, such as great circles on a sphere and closed geodesics on a torus; they satisfy only "locally shortest" or "first variation zero", not necessarily global minimality.}
\ASK{Gauss--Bonnet 公式的意义是什么？}{What is the significance of the Gauss--Bonnet formula?}{它把局部量（Gauss 曲率的积分）与整体拓扑量（Euler 示性数）联系起来：总曲率只由曲面的亏格决定，与具体形状无关。它是绝妙定理的全局版本，也给出测地三角形角度与曲率的关系、多面体角亏损定理、以及平行移动和乐的几何解释。}{It links a local quantity (the integral of Gaussian curvature) with a global topological quantity (the Euler characteristic): the total curvature is determined by the genus of the surface and is independent of its shape. It is the global version of the Theorema Egregium and yields the relation between angles and curvature for geodesic triangles, the angle-defect theorem for polyhedra, and the geometric explanation of holonomy in parallel transport.}

\section*{应用与延伸 | Applications and outlook}
\APS{曲线曲面造型与 CAD。工程中的曲线常用 Bézier 曲线和 B 样条表示，它们由控制多边形生成、具有凸包性与局部修改性；曲面用张量积 Bézier 曲面或 NURBS 表示。几何连续性（$G^1$ 切向连续、$G^2$ 曲率连续）比参数连续性更符合设计需求，而曲率、挠率、主曲率正是判断曲面光顺性的指标。可展曲面（$K=0$）在造船、钣金与建筑幕墙中特别重要，因为只有可展曲面才能由平板材料无拉伸地成形。}{Curve and surface modelling and CAD. Engineering curves are usually represented by Bézier curves and B-splines, generated from control polygons with the convex-hull and local-modification properties; surfaces use tensor-product Bézier surfaces or NURBS. Geometric continuity ($G^1$ tangent continuity, $G^2$ curvature continuity) matches design requirements better than parametric continuity, and curvature, torsion and principal curvatures are exactly the quantities used to judge surface fairness. Developable surfaces ($K=0$) are especially important in shipbuilding, sheet metal and architectural facades, since only developable surfaces can be formed from flat sheets without stretching.}
\APS{地图投影。把球面（地球）画到平面上必然产生失真。Mercator 投影是保角映射，适合航海，因为航向线（等角航线）被画成直线，但高纬度面积被严重放大；等积投影保持面积而牺牲角度；两者不能兼得，因为绝妙定理告诉我们球面 $K=1/R^2\ne0$ 与平面 $K=0$ 不可能等距。这个例子是"内蕴曲率是障碍"的最著名应用。}{Map projections. Drawing a sphere (the Earth) on a plane necessarily distorts it. The Mercator projection is conformal and suits navigation because lines of constant bearing (rhumb lines) become straight lines, but areas near the poles are greatly exaggerated; equal-area projections preserve area at the expense of angles. One cannot have both, because the Theorema Egregium tells us that a sphere with $K=1/R^2\ne0$ cannot be isometric to a plane with $K=0$. This is the most famous application of "intrinsic curvature as an obstruction".}
\APS{极小曲面与肥皂膜。肥皂膜在表面张力下使面积局部极小，其平均曲率 $H=0$，因此悬链面、螺旋面等极小曲面是肥皂膜的自然形状；Plateau 问题研究给定边界下极小曲面的存在性。极小曲面在建筑（张拉膜结构、慕尼黑奥林匹克体育场屋顶）与材料科学（嵌段共聚物、液晶）中都有应用，也是几何分析的核心对象之一。}{Minimal surfaces and soap films. A soap film minimizes area locally under surface tension, so its mean curvature vanishes, $H=0$; hence the catenoid, the helicoid and other minimal surfaces are the natural shapes of soap films, and the Plateau problem studies existence of minimal surfaces with prescribed boundary. Minimal surfaces appear in architecture (tensile membrane structures, the roof of the Munich Olympic Stadium) and in materials science (block copolymers, liquid crystals), and they are a central object of geometric analysis.}
\APS{广义相对论中的曲率思想。Einstein 的广义相对论把引力解释为时空的弯曲：物质分布决定时空的度量，自由下落的粒子沿测地线运动，潮汐力由 Riemann 曲率张量描述；二维的 Gauss 曲率在四维时空中推广为 Riemann 曲率。第 7 章的活动标架与外微分（结构方程、Cartan 引理）会把本讲义的标架方法推广到高维与抽象流形，是学习现代微分几何的入口。}{Curvature in general relativity. Einstein's general relativity interprets gravity as the curvature of spacetime: the distribution of matter determines the metric, freely falling particles move along geodesics, and tidal forces are described by the Riemann curvature tensor; the two-dimensional Gaussian curvature generalizes to the Riemann curvature in four dimensions. Chapter 7, on moving frames and exterior calculus (structure equations, Cartan's lemma), extends the frame method of these notes to higher dimensions and abstract manifolds and is the gateway to modern differential geometry.}

\section*{中英术语表 | Chinese--English glossary}
\noindent\small
\begin{longtable}{@{}>{\raggedright\arraybackslash}p{3.5cm}>{\raggedright\arraybackslash}p{4.6cm}>{\raggedright\arraybackslash}p{6.6cm}@{}}
\toprule
\textbf{中文} & \textbf{English} & \textbf{简要说明 / Note} \\
\midrule
\endhead
正则参数曲线 & regular parametrized curve & 速度处处非零的可微曲线 \\
弧长 & arc length & $s=\int|\bgamma'|\dd t$，单位速度参数 \\
曲率 & curvature & $\kappa=|\bT'|$，切线方向的转动率 \\
挠率 & torsion & $\tau=-\bB'\cdot\bN$，离开密切平面的速率 \\
Frenet 标架 & Frenet frame & $\{\bT,\bN,\bB\}$ 右手正交标架 \\
密切平面 & osculating plane & 由 $\bT,\bN$ 张成的平面 \\
主法向量 & principal normal & $\bN=\bT'/\kappa$ \\
副法向量 & binormal & $\bB=\bT\times\bN$ \\
曲率半径 & radius of curvature & $1/\kappa$ \\
渐屈线 & evolute & 曲率圆心的轨迹 \\
渐伸线 & involute & 切线为原曲线法线的曲线 \\
相对曲率 & relative (signed) curvature & 平面曲线的带符号曲率 \\
正则参数曲面 & regular parametrized surface & $\br_u\times\br_v\ne0$ 的双参数曲面 \\
切平面 & tangent plane & 由 $\br_u,\br_v$ 张成的平面 \\
第一基本形式 & first fundamental form & $I=E\dd u^2+2F\dd u\dd v+G\dd v^2$ \\
第二基本形式 & second fundamental form & $II=L\dd u^2+2M\dd u\dd v+N\dd v^2$ \\
法曲率 & normal curvature & $\kappa_n=II/I$ \\
主曲率 & principal curvature & Weingarten 映射的特征值 \\
主方向 & principal direction & 主曲率对应的切方向 \\
Weingarten 映射 & Weingarten map (shape operator) & $S=-\dd\bn$，自伴线性变换 \\
Gauss 曲率 & Gaussian curvature & $K=\det S=(LN-M^2)/(EG-F^2)$ \\
平均曲率 & mean curvature & $H=\frac12\operatorname{tr}S$ \\
Dupin 标形 & Dupin indicatrix & $II(\bv,\bv)=\pm1$ 的二次曲线 \\
椭圆点/双曲点/抛物点 & elliptic / hyperbolic / parabolic point & 按 $K>0$、$K<0$、$K=0$ 分类 \\
脐点 & umbilic point & 两个主曲率相等的点 \\
可展曲面 & developable surface & $K=0$，可等距展开成平面 \\
直纹面 & ruled surface & 由一族直线生成的曲面 \\
保长对应 & isometry & 保持第一基本形式的映射 \\
保角对应 & conformal map & $I'=\lambda^2I$ \\
等温坐标 & isothermal coordinates & $E=G=\lambda^2$、$F=0$ 的坐标 \\
Gauss--Codazzi 方程 & Gauss--Codazzi equations & $I,II$ 的相容性条件 \\
Gauss 绝妙定理 & Theorema Egregium & $K$ 只由第一基本形式决定 \\
测地曲率 & geodesic curvature & 曲线在曲面内转弯的速率 \\
测地线 & geodesic & $\kappa_g\equiv0$ 的曲线 \\
测地坐标系 & geodesic coordinates & $I=\dd r^2+G\dd\theta^2$ \\
平行移动 & parallel transport & $\nabla_{\bgamma'}V=0$ 的向量场 \\
和乐 & holonomy & 绕闭曲线平行移动的转角 \\
常曲率曲面 & surface of constant curvature & $K$ 为常数，Minding 定理 \\
抽象曲面 & abstract surface & 带度量的二维流形 \\
Gauss--Bonnet 公式 & Gauss--Bonnet formula & $\iint K\dd A=2\pi\chi(M)$ \\
Euler 示性数 & Euler characteristic & $\chi=V-E+F=2-2g$ \\
极小曲面 & minimal surface & $H\equiv0$ 的曲面 \\
内蕴几何 & intrinsic geometry & 只依赖第一基本形式的几何 \\
\bottomrule
\end{longtable}
\normalsize

\section*{参考文献 | References}
\noindent 陈维桓．微分几何．北京：北京大学出版社．第 1 章（预备知识）、第 2 章（曲线论）、第 3 章（曲面的第一基本形式）、第 4 章（曲面的第二基本形式）、第 5 章（曲面论基本定理）、第 6 章（测地曲率和测地线）．\par
\noindent 梅向明，黄敬之．微分几何（第四版）．北京：高等教育出版社．第 1--5 章（曲线论、曲面的第一与第二基本形式、曲面论基本定理、测地线）．\par
\noindent M. P. do Carmo. \emph{Differential Geometry of Curves and Surfaces}. Englewood Cliffs: Prentice-Hall, 1976. Chapters 1--4 (curves, regular surfaces, the geometry of the Gauss map, the intrinsic geometry of surfaces).\par
\EN{Main source: Chen Weihuan, \emph{Differential Geometry}, Peking University Press, Chapters 1--6. Supplementary: Mei Xiangming and Huang Jingzhi, \emph{Differential Geometry} (4th ed.), Higher Education Press, Chapters 1--5; do Carmo, \emph{Differential Geometry of Curves and Surfaces}, Chapters 1--4. References are given at chapter level; page numbers are omitted on purpose.}

\end{document}
